%! TEX TS-program = xelatex
\documentclass[12pt,a4paper,oneside,article]{memoir}%Do not touch this first line ;)
%----------------------------------------------------------------------------------------
%	Metropolia Thesis LaTeX Template
%----------------------------------------------------------------------------------------

%----------------------------------------------------------------------------------------
%	THESIS INFO
%----------------------------------------------------------------------------------------

% All general information (main language, title, author (you), degree programme, major
% option, etc.)
% Edit the file chapters/0info.tex to change these information

% Global information (title of your thesis, your name, degree programme, major, etc.)
\def\bilingual{yes}
\def\thesislang{finnish}
\def\secondlang{english}
\author{Arttu Pennanen}
%License
\def\thesiscopy{by}

%Finnish section, for title/abstract
\def\otsikko{Pragmaattisen funktionaalisen \\ohjelmoinnin arviointi}
\def\tutkinto{Insinööri (AMK)}
\def\kohjelma{Tieto- ja viestintätekniikka}
\def\suuntautumis{Ohjelmistotuotanto}
\def\thesisfi{Insinöörityö}
\def\ohjaajat{
    FM Simo Silander
}
\def\tiivistelma{
    TODO \newline

    Lorem ipsum dolor sit amet, consectetur adipiscing elit. Nunc interdum suscipit nunc. Sed suscipit vulputate nisi id vehicula. Proin mi sapien, tempor sit amet lacus vel, fringilla pulvinar purus. Duis non tempor elit. Donec lobortis est eget pretium faucibus. Maecenas sed convallis ligula. Nunc orci diam, ornare eget quam nec, tincidunt ultricies lacus. Curabitur in lorem ac elit facilisis malesuada id lacinia justo. Quisque placerat, metus quis suscipit interdum, nisl lectus auctor odio, ut consectetur eros leo ac est.
}
\def\avainsanat{funktionaalinen ohjelmointi, JavaScript, TypeScript, koodin luettavuus, koodin ylläpidettävyys, koodin suorituskyky, kehittäjäkokemus}
\def\aihe{TODO: Lyhyt kuvaus opinnäytetyöstä. Max 255 merkkiä. Funktionaalista koodia jne jne}%for the pdf metadata/properties. If not used, empty it and also the \def\subject.

%English section, for title/abstract
\title{Evaluating the Practical Application of Functional Programming}
\def\metropoliadegree{Bachelor of Engineering} % change to your needs, e.g. "master", etc.
\def\metropoliadegreeprogramme{Information and Communication Technology}
\def\metropoliaspecialisation{Software Engineering}
\def\thesisen{Bachelor’s Thesis} % change to your need, e.g. master's
\def\metropoliainstructors{
    Simo Silander, M. Sc.
}
\def\abstract{
    TODO: Abstract content. To force newline between paragraph in the abstract, you must have both a empty line and the newline command. \newline

    Lorem ipsum dolor sit amet, consectetur adipiscing elit. Nunc interdum suscipit nunc. Sed suscipit vulputate nisi id vehicula. Proin mi sapien, tempor sit amet lacus vel, fringilla pulvinar purus. Duis non tempor elit. Donec lobortis est eget pretium faucibus. Maecenas sed convallis ligula. Nunc orci diam, ornare eget quam nec, tincidunt ultricies lacus. Curabitur in lorem ac elit facilisis malesuada id lacinia justo. Quisque placerat, metus quis suscipit interdum, nisl lectus auctor odio, ut consectetur eros leo ac est.
}
\def\metropoliakeywords{functional programming, JavaScript, TypeScript, code readability, code maintainability, code performance, developer experience}
\def\subject{TODO: Short description of the thesis. Max 255 characters. Its the funksönal programming}%for the pdf metadata/properties. If not used, empty it and also the \def\aihe.


%----------------------------------------------------------------------------------------
%	GLOBAL STYLES
%----------------------------------------------------------------------------------------

% If you need extra package, etc. modify the style/style.tex file.
% If you are using Windows OS, you will need to change default font to Arial in that
% style/style.tex file (or install Liberation Sans font to your system).
% If you are using MacOS or linux, make sure you have Liberation Sans font installed.
% Global style. Normally should not be edited.
% If you use windows OS, eventually change \setmainfont to Arial
% Check around commit https://github.com/panunu/metropolia-thesis-latex/commit/a0c15ac77bab1a52c59c517a18080938e57bf5ef
% to see how the font files were manually added (after downloading them: https://pagure.io/liberation-fonts/ )

%\usepackage[l2tabu, orthodox]{nag}%check for obsolete packages (with outdated nag package?)

%condition e.g. for adding or not space in TOC,...
\usepackage{etoolbox}
\ifdefstring{\bilingual}{yes}{
	\usepackage[\secondlang,\thesislang]{babel}% finnish (or swedish) and english
}{
	\usepackage[\thesislang]{babel}% english only
}
\usepackage{iflang}
\usepackage{amsmath}
\usepackage{amsfonts}%extra mathematical symbols
\usepackage{amssymb}
\usepackage{fontspec}
\usepackage{titlesec}
\usepackage{mathtools}%enhance the appearance of documents containing a lot of mathematics
\usepackage{siunitx}  % Load siunitx package for SI units
\usepackage[version=3]{mhchem}
\usepackage{tikz} % mindmaps, flowcharts, piecharts, examples at http://www.texample.net/tikz/examples/
\usetikzlibrary{shapes.geometric, arrows}

\usepackage{ragged2e}
\IfLanguageName{finnish}{
	\RaggedRight%2021 template, align left and hyphennization for Finnish version
	\setlength{\RaggedRightRightskip}{0pt plus 1fil}%TODO: fix the Overfull/Underfull warnings when \RaggedRight (likely in title and abstact)
}{}
\makeatletter
\let\@gnewline\@raggedtwoe@saved@gnewline% Restore original functionality of \newline
\let\\\@raggedtwoe@savedcr% Restore original functionality of \\
\makeatother

%celsius
\newcommand{\celsius}{\ensuremath{^\circ\mathrm{C}}}
%for compact list
\usepackage{enumitem}
%\usepackage{verbatim}%for block comment
%forcing single line spacing in bibliography
\DisemulatePackage{setspace}
\usepackage{setspace}
%\usepackage{eurosym}%euro symbol
%try to count
\usepackage{totcount}
%insert source code
%require -8bit -shell-escape in the xelatex compile command
%if compiling locally, consider options cachedir=minted,outputdir=~/.tex
\usepackage[newfloat]{minted}
\setminted{tabsize=2,linenos,breaklines,breaksymbolleft={\quad},baselinestretch=1}
\setmintedinline{breaklines,breakanywhere}
\usepackage[singlelinecheck=false]{caption}
%force the width of a table instead of column
\usepackage{tabularx}
\usepackage{booktabs} %why not booktabs? :3

\usepackage{csquotes}% avoid warning with babel

\usepackage{float} % For forced figure location with modifier H (\begin{figure}[H])
\usepackage{wrapfig}

%set date format to D.M.YYYY
\def\pvm{\the\day.\the\month.\the\year}
%set date format to D Month YYYY
\usepackage[en,useregional=false]{datetime2}
\DTMsetup{datesep=.}
\DTMnewdatestyle{dMonthyyyy}
{%
	\renewcommand{\DTMdisplaydate}[4]{%
		\number##3 % day
		~% separator
		\DTMenglishmonthname{##2}% month name
		~% separator
		\number##1% year
	}%
	\renewcommand{\DTMDisplaydate}[4]{%
		\DTMdisplaydate{##4}{##3}{##2}{##1}%
	}%
}
\DTMsetdatestyle{dMonthyyyy}
\date{\today}


\usepackage{xcolor}
\usepackage{tcolorbox}
\tcbuselibrary{listingsutf8} % For including monospace font settings

\newcommand{\key}[1]{%
	\begin{tcolorbox}[
			colback=gray!20, % Light gray background
			colframe=gray!20, % Same as background, effectively invisible
			boxrule=0pt, % No border
			arc=4pt, % Rounded corners
			left=2pt, % Left padding
			right=2pt, % Right padding
			top=2pt, % Top padding
			bottom=2pt, % Bottom padding
			boxsep=0pt, % No separation between text and frame, padding only
			height=12pt,
			valign=center,
			nobeforeafter, % No additional space before and after the box
			tcbox raise base, % Aligns the box with the baseline of the surrounding text
			hbox % Ensures the box is only as wide as its contents
		]
		\texttt{\small #1} % Smaller monospaced font for the key text
	\end{tcolorbox}
}

\newcommand\tn[1]{\textnormal{#1}} %use \tn instead of \textnormal
\newcommand\reaction[1]{\begin{equation}\ce{#1}\end{equation}} %\reaction{} for chemical reactions

%NORMAL TEXT
%all text, title, etc. in the same font: Arial
%NOTE: fontname is case-sensitive
\setmainfont[Scale=0.98]{Liberation Sans}
%line space
\linespread{1.46}
\AtBeginEnvironment{tabular}{\singlespacing}
%\doublespacing
%margin
%geometry moved after hyperref
\setlength{\parindent}{0pt} %first line of paragraph not indented
\setlength{\parskip}{16.5pt} %one empty line to separate paragraph
%list with small line space separation
\tightlists
\setlist[itemize]{before=\singlespacing,itemsep=6pt,leftmargin=63pt,labelsep=22pt,topsep=1.5pt,partopsep=0pt,after=\vspace{-22pt}\newline}
\setlist[enumerate]{before=\singlespacing,itemsep=6pt,leftmargin=63pt,labelsep=22pt,topsep=1.5pt,partopsep=0pt,after=\vspace{-22pt}\newline}

%IMAGE - FIGURE
%the figures should be placed in the "figures" folder
\graphicspath{{figures/}}
%figure number without chapter (1.1, 1.2, 2.1) to (1, 2, 3)
\counterwithout{figure}{chapter}
%border around images
\setlength\fboxsep{0pt}
\setlength\fboxrule{0.5pt}
%space after figure caption (and other float elements)
\setlength{\belowcaptionskip}{-7pt}
\setlength{\intextsep}{16.5pt}%space around floats
\AtBeginEnvironment{table}{\addvspace{16.5pt}}

%TABLE
\counterwithout{table}{chapter}

%SOURCE CODE
\newenvironment{code}
{
	\begin{samepage} % Start keeping content on the same page
		\captionsetup{type=listing}
		}
		{
	\end{samepage} % End keeping content on the same page
}

\IfLanguageName{finnish}
{
	\SetupFloatingEnvironment{listing}{name=Koodiesimerkki}
}
{}%was Listaus
%\counterwithout{lstlisting}{chapter}
%moved after begin document, otherwise does not compile

%TOC
%change toc title
\IfLanguageName{finnish}{\addto{\captionsfinnish}{\renewcommand*{\contentsname}{Sisällys}}}{}
%remove dots
\renewcommand*{\cftdotsep}{\cftnodots}
%chapter title and page number not in bold
\renewcommand{\cftchapterfont}{\normalfont}
\renewcommand{\cftchapterpagefont}{\normalfont}
%sub section in toc
\setcounter{tocdepth}{2}
%subsection numbered
\setcounter{secnumdepth}{2}
\renewcommand{\tocheadstart}{}
\renewcommand{\printtoctitle}[1]{\fontsize{13.5pt}{13.5pt}\bfseries #1\vspace*{-4pt}}
%\renewcommand{\aftertoctitle}{\vspace*{-22pt}\afterchaptertitle}
%spacing after a chapter in toc
\preto\section{%
  \ifnum\value{section}=0\addtocontents{toc}{\vskip11pt}\fi
 }
%spacing after a section in toc
\renewcommand{\cftsectionaftersnumb}{\vspace*{-3pt}}
%spacing after a subsection in toc
\renewcommand{\cftsubsectionaftersnumb}{\vspace*{-1pt}}
%appendix in toc with "Appendix " + num
\IfLanguageName{finnish}{
	\renewcommand*{\cftappendixname}{Liite\space}
	\renewcommand{\appendixtocname}{Liitteet}
}{\renewcommand*{\cftappendixname}{Appendix\space}}
%appendix header
\IfLanguageName{finnish}{\def\appname{Liite\space}}{\def\appname{Appendix\space}}

%TITLES
%chapter title
%\clearforchapter{\clearpage}
\titleformat{\chapter}
{\fontsize{14pt}{14pt}\bfseries\linespread{1}}%\clearpage
{\thechapter}{.5cm}{}
\titlespacing*{\chapter}{0pt}{-.42cm}{.5pt}
\titleformat{\section}
{\fontsize{13.5pt}{13.5pt}\normalfont}
{\thesection}{.5cm}{}
\titlespacing*{\section}{0pt}{9pt}{0pt}
\titleformat{\subsection}
{\fontsize{12.7pt}{12.7pt}\normalfont}
{\thesubsection}{.5cm}{}
\titlespacing*{\subsection}{0pt}{11pt}{0pt}


%QUOTE
\renewenvironment{quote}
{\list{}{\rightmargin=0pt\leftmargin=2.2cm\topsep=-14pt}%
	\item\relax\singlespacing}%\fontsize{10pt}{10pt}
{\vspace{8pt}\endlist}

%BIBLIOGRAPHY
%TODO: toggle for Harward v.s. Vancouver
\usepackage[
	backend=biber,
	bibencoding=utf8,%
	citetracker=true,%
	isbn=false,%
	doi=false,%
	url=true,%
	usetranslator=true,%
	bibstyle=ext-authoryear,%
	citestyle=numeric-comp,%
	sorting=none,%
	sortcites=true,%
	block=none,%
	terseinits=false,%
	giveninits=false,%
	maxbibnames=99,%
]{biblatex}

\setlength{\bibitemsep}{11pt}
\setlength{\biblabelsep}{1cm}%with 1cm horizontal gap

\makeatletter
\RequireBibliographyStyle{numeric}
\makeatother

% set right format
\renewcommand*{\multicitedelim}{;\space}
\renewcommand*{\multinamedelim}{;\space}
\renewcommand*{\finalnamedelim}{~\&\space}
\DeclareFieldFormat{biblabeldate}{#1}
\DeclareDelimFormat[bib]{nameyeardelim}{\addperiod\space}
%\DeclareNameAlias{sortname}{last-first} % deprecated
\DeclareNameAlias{default}{family-given}
\DeclareFieldFormat{labelnumberwidth}{#1} % remove () from label number
\DeclareFieldFormat{title}{#1} % normal font for title
\DeclareFieldFormat{citetitle}{#1}
\DeclareFieldFormat{journaltitle}{#1} % remove underline
\DeclareFieldFormat*{url}{%
	\ifentrytype{online}{\IfLanguageName{finnish}{Verkkoaineisto}{Online}\addperiod\space}{}%
	\ifentrytype{video}{Video\addperiod\space}{}%
	\ifentrytype{podcast}{Podcast\addperiod\space}{}%
	\textless\url{#1}\textgreater\addperiod
} % you can modify how to url looks here

\DeclareFieldFormat{urldate}{%
	\ifentrytype{video}{\space Katsottu\space#1}{%
		\ifentrytype{podcast}{\space Kuunneltu\space#1}{\space Luettu\space#1}}%
}

\DefineBibliographyStrings{finnish}{%
	urlfrom = {},%
	fromjapanese = {japanista},%
	fromenglish = {englannista},%
	fromfinnish = {suomesta},%
	fromgerman = {saksasta},%
}
{
	%new cite command: "Vancouver Short"
	\DeclareCiteCommand{\citeVS}
	{\usebibmacro{prenote}}
	{\usebibmacro{author}, \usebibmacro{title}}
	{\multicitedelim}
	{\usebibmacro{postnote}}

	% new cite command: "Vancouver Short Collection" - necessary when referencing whole collections.
	\DeclareCiteCommand{\citeVSc}
	{\usebibmacro{prenote}}
	{\usebibmacro{editor}, \usebibmacro{title}}
	{\multicitedelim}
	{\usebibmacro{postnote}}
}

\newcommand{\engl}[1]{engl. \textit{#1}.}

\DeclareCiteCommand{\citep}[\mkbibbrackets]
{\usebibmacro{cite:init}%
	\usebibmacro{prenote}}
{\usebibmacro{citeindex}%
	\usebibmacro{cite:comp}}
{}
{\usebibmacro{cite:dump}%
	\usebibmacro{postnote}\addperiod}


\addbibresource{biblio.bib}% for biblatex you need out \printbibliography too

%count the appendices (since the chapter counter is reset after \appendix).
%! require to complie 2 times
\regtotcounter{chapter}

% metadata (title, author, lang,...) for accessibility, etc.
\usepackage{hyperref}
\usepackage{hyperxmp}
\def\isolang{\IfLanguageName{finnish}{fi}{en}} %iso code (based on main language)
\ifdefstring{\thesiscopy}{all}{%
	\def\copyen{Copyright \textcopyright\ \the\year{}, \theauthor}
	\def\copyfi{Kaikki oikeudet pidätetään.  \textcopyright\ \the\year{}, \theauthor}
}{%
	\usepackage[type={CC},modifier={\thesiscopy},version={4.0},]{doclicense}
	\def\copyen{\thetitle\ \textcopyright\ \the\year{} by \theauthor\ is licensed under {\doclicenseLongNameRef}}
	\def\copyfi{\otsikko\ \textcopyright\ \the\year{}, jonka tekijä on \theauthor, on julkaistu {\doclicenseLongNameRef} lisenssillä}
}
\hypersetup{%
	pdfdisplaydoctitle,
	%breaklinks=true,%searching for overfull warnings
	pdfencoding=auto,
	bookmarksdepth=subsection,
	unicode=true,
	keeppdfinfo=true,
	pdflang={\isolang},
	pdfmetalang={\isolang},
	pdftitle={\IfLanguageName{finnish}{\otsikko}{\thetitle}},
	pdfkeywords={\IfLanguageName{finnish}{\avainsanat}{\metropoliakeywords}},
	pdfcopyright={\IfLanguageName{finnish}{\copyfi}{\copyen}},
	pdfsubject={\IfLanguageName{finnish}{\aihe}{\subject}},
}

\ifdefstring{\bilingual}{yes}{%metadata (title and copyright) in multiple language
	\XMPLangAlt{\IfLanguageName{finnish}{en}{fi}}{%
		pdftitle={\IfLanguageName{finnish}{\thetitle}{\otsikko}},
		pdfcopyright={\IfLanguageName{finnish}{\copyen}{\copyfi}},
		pdfsubject={\IfLanguageName{finnish}{\subject}{\aihe}},
	}
}{}
\urlstyle{same}

%moved after hyperred as can cause conflicts
\usepackage{pdfcomment}%try the alt text for graphics
\usepackage{accsupp}
\newcommand{\AltText}[2]{\BeginAccSupp{method=pdfstringdef,unicode,Alt={{#1}}}\pdftooltip{{#2}}{{#1}}\EndAccSupp{}}

\usepackage[top=2.5cm, bottom=3cm, left=4cm, right=2cm, nofoot]{geometry}

\usepackage{pgfplots} %simple plots etc
\pgfplotsset{compat=1.16}
\usepackage{pgfplotstable}

% Abbreviations, acronym and glossary, in case of bug, remove temporary the noredefwarn
\usepackage[acronym,toc,nonumberlist,section=chapter,noredefwarn]{glossaries}%xindy,%toc, ,nomain
\newglossarystyle{mystyle}{%
	\setglossarystyle{list}% base this style on the list style
	\renewcommand*{\glossentry}[2]{%
		\item[\glsentryitem{##1}%
		            \glstarget{##1}{\glossentryname{##1}:}]
		\glossentrydesc{##1}\glspostdescription\space ##2}
}
\setglossarystyle{mystyle}

\renewcommand*{\glsclearpage}{}

% Normally, you do not need to modify the title style. It's content comes from the
% chapters/0info.tex file.
% TITLE PAGE
% Normally, you should not edit this file.

\makeatletter
\renewcommand{\maketitle}{
\newgeometry{left=4cm,top=2.7cm}
\thispagestyle{empty}

\includegraphics[width=9.1cm,keepaspectratio]{Metropolia_L_oranssi}

\IfLanguageName{finnish}{
  \vspace*{1.36cm}
}{
  \vspace*{1.18cm}
}
\tn{\large\theauthor\\[9.2mm]{\fontsize{27.5pt}{27.5pt}\normalfont\color[RGB]{155,50,35}\IfLanguageName{finnish}{\otsikko}{\thetitle}}}%\\[22pt]\LARGE\alaotsikko\\[1.75cm]}

\vspace*{2.38cm}
{\setstretch{1.512}
  \IfLanguageName{finnish}{
    Metropolia Ammattikorkeakoulu\\
    \tutkinto \\
    \kohjelma \\
    \thesisfi\\
    \pvm % D.M.YYYY date format for Finnish.
  } {
    Metropolia University of Applied Sciences\\
    \metropoliadegree \\
    \metropoliadegreeprogramme \\
    \thesisen\\
    \@date %D Month YYYY for English
  }
}
\clearpage
\restoregeometry
}
\makeatother



%----------------------------------------------------------------------------------------
%	ABBREVIATION AND GLOSSARY
%----------------------------------------------------------------------------------------

% Add/edit all your acronyms, abbreviations, glossary entries, etc. definitions in
% chapters/0abbr.tex file.
% You can have as many as you wish. Only the ones you use in your text (inserted with
% \gls{} command) will print in the Glossary/Lyhenteet.
% Generate the glossary
\makeglossaries

% Acronyms, abbreviations, etc.


% \newacronym{fp}{fp}{Functionaalinen ohjelmointi}

% Glossary entries


\newglossaryentry{ts}{
	name={TypeScript},
	description={ohjelmointikieli (ts)}
}

\newglossaryentry{js}{
	name={JavaScript},
	description={ohjelmointikieli (js)}
}

\newglossaryentry{functional_programming}{
	name={Funktionaalinen ohjelmointi},
	description={Ohjelmointiparadigma, joka korostaa laskennan mallintamista funktioiden avulla}
}

\newglossaryentry{programming_paradigm}{
	name={Ohjelmointiparadigma},
	description={Ohjelmoinnin lähestymistapa tai malli, joka määrittelee, miten ohjelmat rakennetaan ja miten ohjelmointiongelmia ratkaistaan. Esimerkkejä tunnetuista ohjelmointiparadigmoista ovat funktionaalinen, olio- ja imperatiivinen ohjelmointi}
}

\newglossaryentry{imperative_programming}{
	name={Imperatiivinen ohjelmointi},
	description={Ohjelmointiparadigma, jossa ohjelmoija antaa tarkat ohjeet siitä, kuinka tehtävä suoritetaan vaihe vaiheelta. Tämä lähestymistapa keskittyy ohjelman tilan hallintaan ja muuttamiseen käskyjen avulla, kuten muuttujien asettaminen ja silmukoiden käyttö}
}

\newglossaryentry{declarative_programming}{
	name={Deklaratiivinen ohjelmointi},
	description={Ohjelmointiparadigma, jossa ohjelmoija määrittelee, mitä lopputuloksen tulisi olla, mutta ei yksityiskohtaisesti sitä, miten tämä tulos saavutetaan}
}

\newglossaryentry{object_oriented_programming}{
	name={Olio-ohjelmointi},
	description={Ohjelmointiparadigma, joka perustuu olioiden, eli ohjelmakomponenttien, ympärille, jotka yhdistävät sekä dataa että siihen liittyviä toimintoja. Olio-ohjelmointi korostaa perinnän, kapseloinnin ja polymorfismin kaltaisia periaatteita}
}

\newglossaryentry{set_theory}{
	name={Joukko-oppi},
	description={\protect\engl{set theory} Joukko-oppi on joukkojen ominaisuuksiin perehtynyt matematiikan osa-alue. Joukko-oppi on perustavanlaatuisessa merkityksessä tietorakenteissa. Joukko-opissa olennaisimpia laskutoimituksia ovat esimerksiksi unioni, leikkaus ja joukkoerotus}
}

\newglossaryentry{category_theory}{
	name={Kategoriateoria},
	text={kategoriateoria},
	description={\protect\engl{category theory} Kategoriateoria on matematiikan osa-alue, joka tutkii erittäin yleistävällä tasolla rakenteiden välisiä suhteita ja yhteyksiä. Kategoriateoriaa käytetään funktionaalisen ohjelmoinnin apukeinona}
}

\newglossaryentry{monad}{
	name={Monadi},
	plural={Monadit},
	description={\protect\engl{monad} Monadi on abstrakti tietotyyppi, jota käytetään funktionaalisessa ohjelmoinnissa.}
}

\newglossaryentry{language_agnostic}{
	name={Kieliagnostinen},
	description={\protect\engl{language agnostic} Kieliagnostisuudella tarkoitetaan lähestymistapaa, jossa ei keskitytä tiettyyn ohjelmointikieleen, vaan tarkastellaan peruskäsitteitä ja -periaatteita yleisellä tasolla, sovellettavissa eri kieliin riippumatta niiden syntaksista tai erityispiirteistä}
}

\newglossaryentry{correctness}{
	name={Koodin oikeellisuus},
	description={\protect\engl{correctness} Koodin oikeellisuus on ohjelman kyky täyttää sille asetetut vaatimukset virheettömästi. Toisaalta ilman vaatimuksia koodi ei voi olla oikeellista}
}

\newglossaryentry{pure_function}{
	name={Puhdas funktio},
	description={\protect\engl{pure function} Funktiot, jotka eivät vaikuta ohjelman tilaan ja palauttavat aina saman tuloksen samoilla syötteillä}
}

\newglossaryentry{immutable_data}{
	name={Muuttumaton data},
	description={\protect\engl{immutable data} Data, jota ei voi muuttaa sen luomisen jälkeen}
}

\newglossaryentry{higher_order_function}{
	name={Korkeamman asteen funktio},
	plural={Korkeamman asteen funktiot},
	description={\protect\engl{higher order function} Funktiot, jotka voivat ottaa toisia funktioita syötteenään tai palauttavat funktioita tuloksena}
}

\newglossaryentry{combinator}{
	name={Kombinaattori},
	description={\protect\engl{combinator} Funktionaalisen ohjelmoinnin rakennuspalikka, joka yhdistää funktioita ilman muuttujien käyttöä},
	plural={Kombinaattorit}
}

\newglossaryentry{composed_function}{
	name={Yhdistetty funktio},
	plural={Yhdistetyt funktiot},
	description={\protect\engl{function composition} Useammasta funktiosta koostettu funktio, jossa funktiot suoritetaan toinen toisensa jälkeen niin, että edellisen tuloste anneetaan seuraavan syötteeksi}
}

%----------------------------------------------------------------------------------------
%	DOCUMENT STARTS HERE...
%----------------------------------------------------------------------------------------

\begin{document}
\IfLanguageName{finnish}{
}{
  \raggedright%2021 template, align left, no hyphennization for English version
}
\counterwithout{listing}{chapter}

%----------------------------------------------------------------------------------------
%	TITLE PAGE
%----------------------------------------------------------------------------------------

\input{style/title_headers.tex}
\maketitle
\newpage

%----------------------------------------------------------------------------------------
%	ABSTRACT / Tiivistelmä
%----------------------------------------------------------------------------------------

% If you are international student writing in English, ignore the Finnish abstract.
% If you are Finnish citizen, you must have 2 abstracts, one in Finnish (or Swedish
% depending on your mother tongue) and one in English regardless of the main language of
% your thesis. Normally, you do not need to modify the abstract style. It's content comes
% from the chapters/0info.tex file.
\ifdefstring{\bilingual}{no}{%
  \input{style/abstract_en.tex}
}{%
  \IfLanguageName{finnish}{%order of abstracts based on main language and spacing hell
    % Abstract in Finnish
% Normally, you should not edit this file. Everything comes from chapter/0info.tex

\pagestyle{empty} %remove page number
\newgeometry{top=1.5cm,left=4cm,right=.5cm}
\begin{otherlanguage}{finnish}
{\large\textbf{Tiivistelmä}}
{\renewcommand{\arraystretch}{1.1}
  \begin{tabular}{@{}p{4.7cm} >{\raggedright\arraybackslash}p{10.8cm}@{}}
  Tekijä: & \makeatletter\@author\makeatother
  \\
  Otsikko: & \otsikko
  \\
  Sivumäärä: & \pageref*{LastPage} sivua %+ \total{chapter} liitettä
  \\
  Aika: & \pvm
  \\[6.5mm]
  Tutkinto: & \tutkinto
  \\
  Tutkinto-ohjelma: & \kohjelma
  \\
  Ammatillinen pääaine: & \suuntautumis
  \\
  Ohjaajat: & \ohjaajat
  \\[11mm]
  \cmidrule[.7pt](l{-.15em}r{5.5em}){1-2}
  \multicolumn{2}{>{\raggedright}p{15.5cm}}{
  \vspace{.2mm}
  \tiivistelma
  }
  \\
  \\
  Avainsanat: & \avainsanat
  \\
\end{tabular}
}
\end{otherlanguage}
\restoregeometry
\clearpage


    % Abstract in English
% Normally, you should not edit this file. Everything comes from chapter/0info.tex

\pagestyle{empty} %remove page number
\newgeometry{top=1.77cm,left=4cm,right=.5cm}
\begin{otherlanguage}{english}
  {\large\textbf{Abstract}}
  {\renewcommand{\arraystretch}{1.1}
    \begin{tabular}{@{}p{4.7cm} >{\raggedright\arraybackslash}p{10.8cm}@{}}
      Author:             & \makeatletter\@author\makeatother
      \\
      Title:              & \makeatletter\@title\makeatother
      \\
      Number of Pages:    & \pageref*{LastPage} pages         %+ \total{chapter} appendices
      \\
      Date:               & \makeatletter\thedate\makeatother
      \\[6.5mm]
      Degree:             & \metropoliadegree
      \\
      Degree Programme:   & \metropoliadegreeprogramme
      \\
      Professional Major: & \metropoliaspecialisation
      \\
      Supervisors:        & \metropoliainstructors
      \\[9mm]
      \cmidrule[.7pt](l{-.15em}r{5.5em}){1-2}
      \\
      \multicolumn{2}{>{\raggedright}p{15.5cm}}{
      \makebox[0pt][l]{\hspace{-0.2cm}\parbox{\dimexpr\textwidth-0.5cm\relax}{\abstract}}
      }
      \\
      \\
      Keywords:           & \metropoliakeywords
      \\
    \end{tabular}
  }
\end{otherlanguage}
\restoregeometry
\clearpage


  }{
    \input{style/abstract_en.tex}
    % Abstract in Finnish
% Normally, you should not edit this file. Everything comes from chapter/0info.tex

\pagestyle{empty} %remove page number
\newgeometry{top=2.04cm,left=4cm,right=.5cm}
\begin{otherlanguage}{finnish}
{\large\textbf{Tiivistelmä}}
{\renewcommand{\arraystretch}{1.1}
  \begin{tabular}{@{}p{4.7cm} >{\raggedright\arraybackslash}p{10.8cm}@{}}
  Tekijä: & \makeatletter\@author\makeatother
  \\
  Otsikko: & \otsikko
  \\
  Sivumäärä: & \pageref*{LastPage} sivua %+ \total{chapter} liitettä
  \\
  Aika: & \pvm
  \\[6.5mm]
  Tutkinto: & \tutkinto
  \\
  Tutkinto-ohjelma: & \kohjelma
  \\
  Ammatillinen pääaine: & \suuntautumis
  \\
  Ohjaajat: & \ohjaajat
  \\[4mm]
  \cmidrule[.7pt](l{-.15em}r{5.5em}){1-2}
  \\
  \multicolumn{2}{>{\raggedright}p{15.5cm}}{
  \tiivistelma
  }
  \\
  \\
  Avainsanat: & \avainsanat
  \\
\end{tabular}
}
\end{otherlanguage}
\restoregeometry
\clearpage


  }
}
%----------------------------------------------------------------------------------------
%	License? Acknowledgement?
%----------------------------------------------------------------------------------------

% Uncomment next line and edit chapters/0license.tex if you want license in your thesis.
% License of your thesis
% If you wish to explain what it means. When you publish your thesis in https://theseus.fi
% you will be able to choose between some Creative Commons licenses
% https://creativecommons.org
% Adapt this example text to your taste.
% This would also be the right place to explain the license you choose for the code you
% produced for your thesis.

\pagestyle{empty}
\newgeometry{top=2.04cm,left=4cm}

{\large\textbf{Lisenssit}}

\ifdefstring{\thesiscopy}{all}{}{%
  \begin{wrapfigure}{r}{0.3\textwidth}
    \doclicenseImage
  \end{wrapfigure}
}
\IfLanguageName{finnish}{\copyfi}{\copyen}


\textbf{Voit vapaasti:}
\begin{itemize}
  \item Jakaa {\textemdash} kopioida aineistoa ja levittää sitä edelleen missä tahansa välineessä ja muodossa missä tahansa tarkoituksessa, myös kaupallisesti.
  \item Muunnella {\textemdash} remiksata ja muokata aineistoa sekä luoda sen pohjalta uusia aineistoja missä tahansa tarkoituksessa, myös kaupallisesti.
\end{itemize}

\textbf{Seuraavilla ehdoilla:}
\begin{itemize}
  \item Nimeä {\textemdash} Sinun on mainittava lähde asianmukaisesti, tarjottava linkki lisenssiin sekä merkittävä, mikäli olet tehnyt muutoksia.\\ Voit tehdä yllä olevan millä tahansa kohtuullisella tavalla, mutta et siten, että annat ymmärtää lisenssinantajan suosittelevan sinua tai teoksen käyttöäsi.
  \item Ei muita rajoituksia {\textemdash} Et voi asettaa sellaisia oikeudellisia ehtoja tai teknisiä estoja, jotka estävät oikeudellisesti muita tekemästä mitään sellaista, minkä lisenssi sallii.
\end{itemize}
\bigskip
{\large\textbf{PDF-tiedoston linkit}}

Navigoinnin helpottamiseksi PDF-tiedostossa on käytössä linkkejä. PDF-katseluohjelmasta riippuen linkeissä näkyy ääriviivat:

\begin{itemize}
  \item\fcolorbox{red}{white}{Punaiset ääriviivat}: tiedoston sisäiset linkit sanastoon
  \item\fcolorbox{green}{white}{Vihreät ääriviivat}: tiedoston sisäiset linkit lähdeluetteloon
  \item \fcolorbox{cyan}{white}{Siniset ääriviivat}: ulkoiset linkit verkkosivuille.
\end{itemize}

\noindent Sisäisistä linkeistä voi palata takaisin näppäinyhdistelmällä \key{Alt}+ \key{Vasen nuoli}
(tai \key{Cmd}+ \key{Vasen nuoli}).
\restoregeometry
\clearpage


% Uncomment next line and edit chapters/0acknowledgement.tex if you want acknowledgements.
%\input{chapters/0acknowledgement.tex}

%----------------------------------------------------------------------------------------
%	TABLE OF CONTENTS
%----------------------------------------------------------------------------------------

% Table of Content
% Normally, you don't have to modify this file.
% In case spacing of appendix/liite goes crazy, look in ../thesis.tex and look
% for the hack described in APPENDICES.

\makeevenhead{plain}{}{}{}
\makeoddhead{plain}{}{}{}
\pagestyle{empty} %remove page number in toc (if longer than 2 pages)
\newgeometry{top=2.04cm, left=4cm}
\tableofcontents*
\pagestyle{empty} %remove page number in toc (if longer than 1 pages)

\restoregeometry
\clearpage



%list of figure, tables would come here if relevant?

%----------------------------------------------------------------------------------------
%	Lyhenteet / Abbreviation
%----------------------------------------------------------------------------------------

% If you don't use abbreviations/glossary, remove the following line.
% Abbreviation and Glossary
% Normally, you don't have to modify this file. Your abbreviations, etc. goes in
% ../chapters/0abbr.tex file.

\begin{singlespacing}

	% \gsladdall%would add all terms even if not used in your text.
	\addtocontents{toc}{\cftpagenumbersoff{chapter}}
	{
		\titleformat{\section}
		{\fontsize{13.5pt}{13.5pt}\bfseries\normalfont}
		{\thesection}{.5cm}{}
		\renewcommand*{\glossarypreamble}{\vspace{.7\baselineskip}}
		%Adapt labelwidth (sorry for the ugly hack)
		\setlist[description]{leftmargin=!, labelwidth=4em, before=\doublespacing}
		\IfLanguageName {finnish} {
			\printacronyms[title=Lyhenteet, nogroupskip]
		}{
			\printacronyms[title=List of Abbreviations, nogroupskip]
		}
		\setlist[description]{leftmargin=!, labelwidth=7em, before=\doublespacing}
		\vspace{1cm}
		\printglossary[nogroupskip]
		\setlist[description]{style=standard} % reset settings back to default
	}
	\addtocontents{toc}{\cftpagenumberson{chapter}}
\end{singlespacing}

\clearpage


%----------------------------------------------------------------------------------------
%	CONTENT
%----------------------------------------------------------------------------------------

\input{style/content.tex}%reset page number to 1, etc.

% Introduction

\chapter{Johdanto}


Insinöörityössä tutkitaan funktionaalisen ohjelmoinnin käytänteitä sellaisissa ohjelmointikielissä, joissa funktionaalinen ohjelmointi ei ole ollut kielen luonnin peruste. Etsitään pragmaattista, tavoitelähtöistä puolta funktionaalisesta ohjelmoinnista, jossa teoriaa ei painoteta ilman selkeitä ja perusteltuja syitä. Funktionaalinen ohjelmointi on perusteiltaan erittäin teoreettista ja akateemisista lähtökohdista.

Kokemusten ja median tarkkailun perusteella on huomattu funktionaalisten periaatteiden yhdistyvän yleiseen näkemykseen siitä, mitä ohjelmoinnin tulisi olla. Funktionaalinen ohjelmointi ei ole kuitenkaan noussut valtavirran suosioon. Insinöörityön motivaationa on etsiä toimintaperiaatteita funktionaalisen ohjelmoinnin tuontiin mukaan arkiseen ohjelmointiin ilman pakonomaista teoriaa. Insinöörityössä kuitenkin etsitään funktionaalisen ohjelmoinnin teoriasta niitä osia, joilla voitaisiin kasvattaa funktionaalisen ohjelmoinnin pragmaattista mukaanottoa.

Insinöörityön koodiesimerkeissä käytetään \glsdisp{ts}{TypeScriptiä} ja \glsdisp{js}{JavaScriptiä}. Vaikka kaikki validi JavaScript koodi on validia TypeScript koodia \cite{typsecript_website}, niin kielten nimiä käytetään työssä tilanteeseen sopien. Kirjoitetaan TypeScriptistä, kun sen tarjoamat tyypit JavaScriptin päälle ovat tilanteeseen nähden merkittäviä. Toisaalta kirjoitetaan JavaScriptistä, kun tyypit eivät ole esimerkille merkittäviä. Muuten pyritään pitämään käsitellyt asiat sellaisina, että niitä voisi hyödyntää myös muissa ei-funktionaalisissa ohjelmointikielissä.

% Motivation
\clearpage

\chapter{Motivaatio}

Ei ole yksiselitteistä, mitä \glsdisp{functional_programming}{funktionaalisen ohjelmoinnin} tulisi olla ja miten sitä tulisi käyttää. Matemaatikot vannovat funktionaalisen ohjelmoinnin \glsdisp{correctness}{oikeellisuuden} nimeen, ja pragmaatikot kauhistuvat. Kuuluisan ristiriitaiset ja mutkikkaat keskustelut hämärtävät pohjaa, minkä päälle funktionaalinen ohjelmointi on rakennettu. \citep{stackoverflow_what_monad,stackoverflow:why_monad,promises-spec-94,dear_functional_bros}

Funktionaalisen ohjelmoinnin hyötyjä on usein sanottu löytyvän suurelta osin vain matemaattisissa tai datapainotteisissa ohjelmissa \cite[10]{cantarella_fp_haitat}. Ajatus on haastettavissa.

Ensimmäisenä huomiona on se, että suuri osa ohjelmista on juuri datapainoitteisia \cite{gartnerb2b}. Tietokannassa on dataa, jota näytetään käyttöliittymässä. Samaan aikaan kerätään lisää dataa, jota näytetään toisissa käyttöliittymissä; toisin sanoen varsin datapainoitteista.

Toisena huomiona on se, että kaikki ongelmat voidaan mallintaa matemaattisesti. Funktionaalisen ohjelmoinnin Turing-vahva lambda-kalkyyli-pohja määritelmällisesti mahdollistaa sen, että funktionaalisessa ohjelmoinnissa voi toteuttaa kaiken saman kuin olio-ohjelmoinnissa \cite{Tan2004,computerphile_lambda}.

Ongelmien mallintaminen matemaattisesti ei kuitenkaan ole helppoa. Mallintaminen tuntuu vaativan vahvoja tyyppijärjestelmiä mallinnuksen järkevästi todistamiseksi.

Kuitenkin pikkuhiljaa funktionaalista ohjelmointia aletaan kutsumaan ilmiselväksi reitiksi ohjelmoida, ja esimerkiksi olio-ohjelmointia ylivertaisemmaksi \cite[1]{the_oo_way}. Jopa Robert C. Martin, tunnettu ohjelmointikäytäntöjen asiantuntija ja olio-ohjelmoinnin puolestapuhuja, on siirtynyt merkittävästi funktionaalisen ohjelmoinnin suuntaan \cite{martin2019whyclojure,martin2017pragmaticfp}. Lienee tarpeen selvittää, miten funktionaalisuutta saataisiin enemmän osaksi nykyohjelmoijan työkalupakkia.
\clearpage

Tiukasti funktionaaliseen ohjelmointiin tarkoitetut kielet eivät tunnu saavan nostetta päästäkseen yleiseen suosioon (kuva \ref{fig:fplangpopularity}). Ero siirtyessä \glsdisp{object_oriented_programming}{olio-ohjelmoinnista} tai \glsdisp{imperative_programming}{imperatiivisesta ohjelmoinnista} funktionaaliseen ohjelmointiin voi olla kuin yöllä ja päivällä. Siksi funktionaalisen ohjelmoinnin keinoja ripotellaan osaksi suosittuja ohjelmointikieliä, että saataisiin vietyä todettuja parhaita funktionaalisen ohjelmoinnin paloja myös muihin paradigmoihin. Kenties ajatellaan niin, että varpaat on ensin kasteltava, ennen kuin voi hypätä syvään päätyyn.

\begin{figure}[htbp]
    \pgfplotstableread[col sep=comma]{data/fp_lang_popularity.csv}\datatable
    \pgfplotstablegetrowsof{\datatable}
    \pgfmathtruncatemacro{\rows}{\pgfplotsretval-1}
    \centering
    \begin{tikzpicture}
        \begin{axis}[
                width=\linewidth,
                height=8cm,
                ylabel={Suosio},
                xlabel={Vuosi},
                xtick={0,12,...,\rows},
                xticklabel style={rotate=45, anchor=east},
                xticklabels={2004, 2005, 2006, ..., 2024},
                legend pos=north west,
                grid=major,
                cycle list name=color list
            ]
            \addplot [mark=none, color=blue] table [x expr=\coordindex, y=Erlang] {\datatable};
            \addplot [mark=none, color=red] table [x expr=\coordindex, y=Haskell] {\datatable};
            \addplot [mark=none, color=green] table [x expr=\coordindex, y=OCaml] {\datatable};
            \addplot [mark=none, color=orange] table [x expr=\coordindex, y=Elixir] {\datatable};
            \addplot [mark=none, color=purple] table [x expr=\coordindex, y=Scala] {\datatable};
            \legend{Erlang, Haskell, OCaml, Elixir, Scala}
        \end{axis}
    \end{tikzpicture}
    \caption{Funktionaalisten ohjelmointikielien Erlangin, Haskelin, OCamlin, Elixirin ja Scalan suosiokehitys vuosina 2004\textendash2024. Kielet eivät ole olleet suuressa nousussa. \citep{fplanggoogletrend}}
    \label{fig:fplangpopularity}
\end{figure}


Ripottelun voi havaita esimerkiksi, kun vahvasti olio-ohjelmointiin sovellettua Javaa on tuotu kohti funktionaalista pelikenttää erilaisin kielen työkalujen, kuten lambdafunktioiden ja Stream-API:n metodien yhdistämisen avulla \cite{oracle_function_package,oracle_stream_api}.
Javaa voikin nykyään jo kutsua hybridiksi ohjelmointikieleksi funktionaalisen ja olio-ohjelmoinnin välillä \cite[50]{sundstrom_java_fp}.

Myös JavaScriptissä Array-metodit, kuten \mintinline{javascript}{Array.prototype.map} \cite{mdn_array_map}, \mintinline{javascript}{Array.prototype.filter} \cite{mdn_array_filter} ja \mintinline{javascript}{Array.prototype.reduce} \cite{mdn_array_reduce}, ovat vakiintuneet keskeisiksi työkaluiksi \cite{vakil2016,8_must_know_array_methods}. Kyseiset funktiot ovat laajalti käytössä funktionaalisessa ohjelmoinnissa. Myös \gls{js} Set -tietorakenne (joukko) on juuri saamassa \glsdisp{set_theory}{joukko-opille} ja funktionaaliselle ohjelmoinnille olennaisia metodeja kuten \mintinline{javascript}{Set.prototype.union} \cite{mdn_set_union} ja \mintinline{javascript}{Set.prototype.difference} \cite{mdn_set_intersection}, jotka palauttavat uusia Set-tietorakenteita (joukkoja) aiempien muokkaamisen sijaan funktionaaliselle ohjelmoinnille ominaisesti \cite{mdn_set_methods}.

Myös kirjoitushetkellä tason 3 TC39-ominaisuusesitys \gls{js} iteraattorien apumetodeille vahvistaa funktionaalista tapaa ohjelmoida \cite{tc39_iterator_helpers}. Näiden metodien lisääminen mahdollistaa vielä johdonmukaisemman funktionaalisen koodin kirjoittamista.

\section{Kokemukset funktionaalisesta ohjelmoinnista}

Tehtäessä työtä osana Hoxhunt Oy:n tuotekehitystiimiä on saatu kokemusta funktionaalisesta ohjelmoinnista. Kokemuspohjaista tietoa on kerätty tiimiläisiltä sekä muuten omalta osalta. On nähty, kuinka funktionaalisen ohjelmoinnin käytänteitä käytetään, siten, että ensin niiden käyttöön on kannustettu, ja toisin pyritty estämään. Koodikatselmuksissa ja muissa keskusteluissa on usein kuultu ajatuksia kärkkäästi niin funktionaalisen ohjelmoinnin puolesta kuin vastaankin.

Funktionaalisen ohjelmoinnin kulmakivi on sen matemaattinen perusta \cite{computerphile_lambda,Tan2004}. Toisinaan funktionaalista ohjelmointia on mielletty vain älykkyyden todistamiseksi, ja jota humoristisesti ajatellaan vain matematiikan maisterin tutkinnon suorittaneiden taitavan. Voi olla vaikea nähdä funktionaalisen ohjelmoinnin hyötyjä, jos kokee sen vain ylimääräisenä monimutkaisuutena.

Funktionaalisen koodin kääntöpuoli on sen selittävyys: deklaratiivisuus. Jos on nähty funktionaalisen ohjelmoinnin tuntemattomuuden tai havaitun monimutkaisuuden läpi, niin on huomattu, että funktionaalista ohjelmakoodia voi parhaimmillaan lukea kuin selkeää englannin kieltä. Tämä johtuu siitä, että funktionaalinen ohjelmakoodi on perusteiltaan (\glsdisp{declarative_programming}{deklaratiivista ohjelmakoodia}) suorien käskyjen sijaan (\glsdisp{imperative_programming}{imperatiivinen ohjelmakoodi}) \cite{ms:fp_vs_imperative}.

\glsdisp{composed_function}{Yhdistetyt funktiot}, joissa funktioita suoritetaan toinen toisensa jälkeen, ja annetaan aina edellisen paluuarvo seuraavan syötteeksi, kertoo suoraan, mitä siinä tapahtuu. Koodi dokumentoi itse itseään, kun kaikki ohjelman vaiheet ovat selkeästi nimetyissä palikoissa; erillisiä kommentteja ei koodiin tarvitse edes lisätä  (koodiesimerkki \ref{code:first_pipeline}).

\begin{code}
    \begin{minted}{typescript}
    const congratulateUsers = pipe(
        filter(isActiveUser),
        filter(isBirthdayToday),
        map(createCongratulationMessage),
        sendCongratulations
    )
    // Onnitellaan kaikkia asiakkaita, joilla on tänään syntymäpäivä!
    congratulateUsers(users)
\end{minted}
    \caption{JavaScript-esimerkki yhdistetystä funktiosta. Funktioiden yksityiskohtia ei tarvitse tietää, että voi ymmärtää ohjelman kokonaiskuvan.}

    \label{code:first_pipeline}
\end{code}


Toisaalta koodin deklaratiivisuuden ylistämisellä voi yrittää peittää muita ongelmia. Jos sokeana seuraa sääntöjä ja käytänteitä, voi jäädä muuten tärkeitä osa-alueita huomioimatta \cite{functional_fixedness}. Jos pyrkii liiallisesti luettavuuteen deklaratiivisella koodilla, voi esimerkiksi ohjelman tehokkuus kärsiä.

Luettavuus, tehokkuus ja ylläpidettävyys ovat kolme keskeistä osa-aluetta, joita on painotettu eniten työpaikalla käydyissä keskusteluissa, kun arvioidaan funktionaalisen ohjelmointitavan laajempia hyötyjä. Nämä osa-alueet näkyvät merkittävinä myös muiden ohjelmointityylien arvioinnissa. Luettavuutta pidetään yleensä kaikkein tärkeimpänä. Ohjelmakoodin ei tarvitse olla tehokasta ennen kuin optimointia todella tarvitaan \cite{prematureoptimization}, ja ylläpidettävyyden perusta on vahvasti sidoksissa siihen, kuinka luettavaa koodi on.

\section{Lähestymistapa}

Tässä insinöörityössä tutkitaan, mitä tapahtuu, kun funktionaalista ohjelmointia ripotellaan proseduraaliseen ohjelmakoodiin. Tätä tarkastellaan sellaisen ohjelmointikielen näkökulmasta, jota ei ole suunniteltu erityisesti funktionaaliseen ohjelmointiin, mutta jossa se on kuitenkin mahdollista. Täksi ohjelmointikieleksi valikoitui \gls{js} (ja \gls{ts}). Ohjelmointikieli valittiin aiemman osaamisen ja kokemusten perusteella, ja toisaalta myös siksi, että kieli on yksi markkinoiden käytetyimpiä, ja siten myös yksi olennaisimmista \cite{pypl:lang}. \Gls{js} on myös kielenä siitä mielenkiintoinen, että se mahdollistaa koodiabstraktioiden luomisen ilman suurta vaivaa ohjelmoijalta. Hätiköidyt abstraktiot johtavat epäjohdonmukaisuuteen ja lopulta uudelleenkäyttämättömyyteen \cite{dijkstra_humble_programmer}. Tämä tekee JavaScriptista haastavan, mutta antoisan kielen tutkia pragmaattista funktionaalisen ohjelmoinnin ripottelua.


Insinöörityössä käsiteltäviä asioita pyritään perkaamaan \glsdisp{language_agnostic}{kieliagnostisesti}, jolloin saadut tiedot ja taidot olisivat oleellisia myös kehittäjille, jotka eivät käytä TypeScriptiä tai JavaScriptiä.

Tarkoituksena on löytää tasapaino sille, miten funktionaalinen ohjelmointi kannattaa ottaa mukaan ohjelmistoprojektiin. Pohjalla on ajatus, että paras tapa olisi vain siirtyä kokonaan funktionaalisiin ohjelmointikieliin hybridikielten sijasta. Tätä ajatusta pyritään kumoamaan, sillä toivotaan, että funktionaalisuus voisi toimia myös kevyemmin osana ohjelmoijan arkea.

Kun jotain ei vaadita, se jää usein tekemättä tai ainakin puutteelliseksi. Siksi voi olla haastavaa ohjelmoida pragmaattista funktionaalista koodia JavaScriptissä. Kielessä, jossa juuri mitään sääntöjä ei pohjimmiltaan ole.


% Theoretical background
%\clearpage % Uncomment if needed to force a page break before the chapter
\vspace{21.5pt}
\chapter{Teoreettinen tausta}

Jotta voidaan löytää pragmaattinen näkökulma funktionaaliseen ohjelmointiin, on käsiteltävä funktionaalisen ohjelmoinnin teoreettisia perusteita ja keskeisiä käsitteitä.

Tarkastellaan joukko-opin sovelluksia ohjelmoinnissa käytännön tasolla. Sivutaan kategoriateorian osuutta funktionaalisessa ohjelmoinnissa ilman syvällistä paneutumista yksityiskohtiin. Sivutaan myös \glsdisp{category_theory}{kategoriateoriaa}, joka on funktionaalisen ohjelmoinnin teoreettisen taustan kulmakivi. Tarkemmin käydään läpi kategoriateorian \glsdisp{monad}{monadi}-rakennetta, ja sen toimintaa.

Tutkitaan, miten nämä käsitteet integroituvat moderniin ohjelmistokehitykseen, kuten TypeScriptin käyttöön.

On tärkeää huomata, että tämän insinöörityön tavoitteena ei ole tarjota kattavaa opasta kaikkeen funktionaaliseen ohjelmointiin, vaan keskittyä valittuihin aihealueisiin, jotka tukevat työn pääteemoja ja osoittavat, miten teoreettiset käsitteet voidaan soveltaa käytännön ohjelmointihaasteisiin.

Valitut aihealueet pyritään avaamaan \glsdisp{language_agnostic}{kieliagnostisesti}, mutta kuitenkin JavaScript- sekä TypeScript-ekosysteemien näkökulmasta.

Lopulta näillä tiedoilla voidaan näyttää, miten funktionaalisen ohjelmoinnin käytänteet näkyvät hyödyllisesti (tai haitallisesti) ei-funktionaalisissa ohjelmistoprojekteissa.

Teoriaosuudessa pyritään pitämään sisältö mielenkiintoisena liimaamalla käsitteet yhdeksi kokonaisuudeksi. Teoriaan pyritään tuomaan konkretiaa erilaisilla esimerkeillä ja anekdooteilla.

\section{Funktionaalinen ohjelmointi}

Funktionaalinen ohjelmointi on ohjelmointiparadigma, jonka juuret ulottuvat 1930-luvulle ja lambda-kalkyylin kehitykseen. Lähestymistapa korostaa funktioiden käyttöä perusyksikköinä ja pyrkii välttämään \glsdisp{immutable_data}{muuttuvia tiloja} (mutability) sekä \glsdisp{side_effects}{sivuvaikutuksia} (side-effects). Funktionaalisen ohjelmoinnin keskeisiä käsitteitä ovat \glsdisp{pure_function}{puhtaat funktiot}, \glsdisp{higher_order_function}{korkeamman asteen funktiot}, \glsdisp{composed_function}{yhdistetyt funktiot}, sekä \gls{declarative_programming}. Funktiot ovat funktionaalisessa ohjelmoinnissa tosiaankin keskiössä. \citep{Tan2004,computerphile_lambda}

\subsection{Työn nimeämiskäytänteet}

Ohjelmoidessa funktionaalista ohjelmakoodia nimeämiskäytänteet ovat yhtä vaikeita kuin aina ennenkin. Tässä insinöörityössä käytetään Haskell-ohjelmoijien keskuudessa näkyviä nimeämiskäytänteitä, joissa käytetään lyhyitä ja \textquote{intuitiivisia} nimiä funktioille ja muuttujille. Esimerkiksi:

\begin{itemize}
    \item x: Yleisesti käytetty muuttuja tai parametri, joka edustaa arvoa, jota funktiot käsittelevät. Jos näkyvyysalalla on muita muuttujia, niitä nimetään usein aakkosjärjestyksessä x:stä eteenpäin (x, y, z).
    \item xs: Muuttuja, joka edustaa useita arvoja. Usein sisältää listan arvoista, mutta tietorakenteen tietäminen ei ole merkittävää.
    \item f: Yksittäinen funktio, jota voidaan käyttää käsittelemään arvoja. Jos näkyvyysalalla on muita funktioita, niitä nimetään usein aakkosjärjestyksessä f:stä eteenpäin (f, g, h, i...).
    \item fs: Muuttuja, joka edustaa useita funktioita. Usein sisältää listan funktioista, mutta tietorakenteen tietäminen ei ole merkittävää.
\end{itemize}

Esimerkkejä käytänteiden käytöstä:

\begin{code}
    \begin{minted}{javascript}
const map    = f => xs => xs.map(f)
const either = f => g => x => f(x) || g(x)
const max    = x => y => x > y ? x : y
const pipe   = fs => x => fs.reduce((acc, f) => f(acc), x)
\end{minted}
    \caption{Esimerkkejä insinöörityössä käytettävistä nimeämiskäytänteistä. Muunmuassa funktio \texttt{map} ottaa ensin yhden funktion (f), ja tämän jälkeen monta arvoa (xs). \texttt{Pipe} ottaa ensin monta funktiota (fs), ja tämän jälkeen yhden arvon (x)}
    \label{code:javascript_naming_convention_example}
\end{code}

\subsection{Puhtaat funktiot}

Funktionaalisessa ohjelmoinnissa \glsdisp{pure_function}{puhtaat funktiot} ovat funktioita, jotka poikkeuksetta aina palauttavat samalle syötteelle saman arvon. Funktion puhtauden voi päätellä siitä, jos sen voisi teoreettisesti korvata hakutaulukolla. \citep{feldman_fp_pragmatists}

Puhtaista funktioista hyötyy siten, että niitä voi käyttää uudelleen ja uudelleen kontekstista riippumatta. Puhtailla funktioilla voi rakentaa järjestelmäriippumattomia kirjastoja.

Puhtaiden funktioiden perusta on matematiikassa. Jos matemaattinen funktio $f(x) = x+2$ esimerkiksi ei aina palauttaisi samalle $x$ arvolle samaa palautusarvoa, mitä hyötyä funktiosta edes olisi? Toivottavasti $2 + 2$ tulee aina olemaan $4$.

\subsection{Kombinaattorit}

Funktionaalinen ohjelmointi mahdollistaa koodin uudelleenkäytettävyyden. Lambda-kalkyylin perustein joka ikisien ongelman voi purkaa pieniksi funktioiksi \cite{BlellochHarper2015}. Lambda-kalkyylistä johdetut kombinaattorit voi ottaa osaksi ohjelmoijan työkalupakkia missä tahansa ohjelmointikielessä, joka kohtelee funktioita samalla tavalla kuin muuttujia (first-class citizen) (Koodiesimerkki \ref{code:javascript_combinators}).

Jos kombinaattoreille haluaa etsiä vastaavanlaista määritelmää olio-ohjelmoinnin termeistä, voisi niitä ajatella olevan perustavanlaatuisia suunittelumalleja.

Kombinaattori \mintinline{javascript}|compose| on tapa tehdä yhdistettyjä funktioita, joista kerrotaan enemmän seuraavassa osassa.

\begin{code}
    \begin{minted}{javascript}
const identity      = x => x
const constant      = x => y => x
const apply         = f => x => f (x)
const thrush        = x => f => f (x)
const duplication   = f => x => f (x) (x)
const flip          = f => y => x => f (x) (y)
const compose       = f => g => x => f (g (x)) 
const substitution  = f => g => x => f (x) (g (x))
\end{minted}
    \caption{Yleiset kombinaattorit esitettynä JavaScriptissä \cite{javascript_combinators}. Kombinaattoreilla voi esittää lambda-kalkyyliä. Kombinaattoreilla voi muokata funktioiden toimintaa, ja yhdistää funktioita yhdiste funktioiksi}
    \label{code:javascript_combinators}
\end{code}

Kombinaattorit näyttävät uhkaavilta, ja raa'assa muodossaan niiden hyötyarvoa voi olla hankala havaita. Niiden käyttämistä voi kuitenkin perustella niiden perustavanlaatuisten ja kieliagnostisien ominaisuuksien takia. Ohjelmoidessa funktionaalisesti ei kombinaattoreita kuitenkaan tarvitse paljoa ajatella. Niitä tulee kirjoittamaan varmasti vähintään vahingossa, sillä niitä vaaditaan funktioita pyörittäessä \cite{javascript_combinators}. Kuitenkin on hyvä tiedostaa käytössä olevat rakennuspalikat, ja millä nimellä etsiä niistä tietoa tarvittaessa.

\subsection{Yhdistetyt funktiot ja niiden vahva merkitys}

Usein funktionaalista ohjelmointia mainostaessa puhutaan funktioiden puhtauden ja datan muuttumattomuuden olevan paradigman oleellisinta. Kuitenkin samaan aikaan juuri datan muuttumattomuus ja tilattomuus koetaan paradigman suureksi kompastuskiveksi \cite{cantarella_fp_haitat,is_reduce_bad,vakil2016}.

On tärkeää huomata, että kaiken ei tarvitse olla täysin puhdasta ja muuttumatonta, kunhan ne osat, joissa funktionaalista ohjelmointia hyödynnetään, pysyvät ennustettavina ja hallittavina.

\glsdisp{composed_function}{Yhdistetyt funktiot} ovat loistava tapa kirjoittaa selkeää ja modulaarista koodia. Parasta on, että kaiken ei tarvitse olla pelkkiä funktioita: voi hyvin kirjoittaa myös olio-ohjelmointityylillä ja silti hyödyntää yhdistettyjä funktioita logiikan selkeyttämiseksi ja yksinkertaistamiseksi.

Loppujen lopuksi funktionaalisen ohjelmoinnin perusta on funktioissa ja niiden yhdistämisessä. Datan muuttumattomuus ja sivuvaikutuksettomuus on vain esivaatimus funktioiden ennustettavuudelle.

\glsdisp{composed_function}{Yhdistetty funktio} tarkoittaa yksinkertaisesti kahden, tai useamman, funktion yhdistämistä siten, että yhden funktion tulos syötetään seuraavalle. Esimerkiksi koodiesimerkin \ref{code:javascript_manual_composition} funktio $h$ on funktioiden $f$ ja $g$ yhdiste. Ajaessa funktio $h$, suoritetaan ensin $g$, jonka palautusarvo annetaan funktiolle $f$. Kaksi funktiota yhdistämällä on saatu yksi uusi funktio.

\begin{code}
    \begin{minted}{javascript}
const f = x => 2 * x

const g = x => x + 3

const h = x => f(g(x))
\end{minted}
    \caption{JavaScript-esimerkki yhdistetystä funktiosta h ilman pipe tai compose funktiota}
    \label{code:javascript_manual_composition}
\end{code}

Funktionaalisissa ohjelmointikielissä yhdistettyjä funktioita pystyy usein kirjoittamaan käyttäen kieleen sisäänrakennettuja operaattoreja, joilla funktioiden yhdistäminen on helppoa ja suoraan osana ohjelmointikieltä \cite{fsharpcomposition,haskellcomposition}.
JavaScriptissä ei ole vastaavanlaista sisäänrakennettua operaattoria.
(Vaikkakin sellainen on ollut kehitteillä \cite{tc39_pipeline_operator}.)
Operaattorin voi kuitenkin kirjoittaa funktiona helposti osaksi mitä tahansa koodikantaa (Koodiesimerkki \ref{code:javascript_pipe_composition}).

\begin{code}
    \begin{minted}{javascript}
const pipe    = f => g => x => g(f(x))
const compose = f => g => x => f(g(x))

const h = pipe(f)(g)
// tai vaihtoehtoisesti
const h = compose(g)(f)
\end{minted}
    \caption{JavaScript-esimerkki funktiokompositiosta pipe ja compose funktioilla}
    \label{code:javascript_pipe_composition}
\end{code}

Voi olla mielenkiinoista huomata, että koodiesimerkin \mintinline{javascript}{compose} löytyy myös koodiesimerkin \ref{code:javascript_combinators} tunnetuista yleisistä kombinaattoreista.

Koodiesimerkissä \ref{code:javascript_pipe_composition} on näytettynä kaksi eri tapaa yhdistää funktioita: \texttt{pipe} ja \texttt{compose}. Käytännössä tapojen ainoa ero on se, kummasta suunnasta funktiot suoritetaan. \texttt{compose} on lähempänä yhdistettyjen funktioiden matemaattisia perusteita, kun taas \texttt{pipe} on usein helpompi lukea, sillä se suoritetaan yleisessä lukemissuunnassa, eli vasemmalta oikealle, tai ylhäältä alas. \citep{whyprefercompose}

Tulevissa koodiesimerkeissä tullaan suosimaan funktioiden yhdistämistä käyttäen \texttt{pipe}-funktiota \texttt{compose}-funktion sijasta.

Esimerkin \ref{code:javascript_pipe_composition} \texttt{pipe} ja \texttt{compose} toimivat vain kahdelle funktiolle kerrallaan. Tämä tapa seuraa lambda-kalkyyliä \cite{computerphile_lambda}. Ohjelmointikielen mukaan voi kuitenkin olla mieluisaa kirjoittaa funktiot niin, että ne tukevat suoraan mielivaltaista määrää funktioita.

JavaScriptissä näin voi tehdä esimerkiksi hyödyntämällä Array-tietorakennetta (Koodiesimerkki \ref{code:javascript_better_pipe}).

\begin{code}
    \begin{minted}{javascript}
const pipe    = fs => x => fs.reduce((acc, f) => f(acc), x)
const compose = fs => x => fs.reduceRight((acc, f) => f(acc), x)

const f = x => x + 1
const g = x => x * 2
const h = x => x - 3

const i = pipe([f, g, h])
i(5) // (((5 + 1) * 2) - 3) = 9
// tai vaihtoehtoisesti
const i = compose([h, g, f])
i(5) // (((5 + 1) * 2) - 3) = 9
\end{minted}
    \caption{JavaScript-esimerkki yhdistettyjen funktioiden luomisesta käyttäen \texttt{reduce} ja \texttt{reduceRight} funktioita. Nämä versiot voivat ottaa mielivaltaisen määrän funktioita argumenttina}
    \label{code:javascript_better_pipe}
\end{code}

Mikä tästä sitten tekee loistavaa? Se, että ohjelmointi on kerrankin oikeasti kuin LEGO palikoilla leikkimistä (Koodiesimerkki \ref{code:javascript_composition_example}).

\begin{code}
    \begin{minted}{javascript}
const pipe     = fs => x => fs.reduce((acc, func) => func(acc), x)
const multiply = x => y => x * y
const add      = x => y => x + y
const filter   = f => xs => xs.filter(f)
const map      = f => xs => xs.map(f)
const isEven   = x => x % 2 === 0

const rejectOdds   = filter(isEven)
const multiplyBy10 = multiply(10)

const pipeline = pipe([rejectOdds, map(multiplyBy10)])

pipeline([1, 2, 3, 4]) // [20, 40]
\end{minted}
    \caption{JavaScript-esimerkki yhdistettyjen funktioiden käyttämisestä laskutoimituksiin. Erilaiset operaatiot on pilkottu uudelleenkäytettäviksi funktioiksi}
    \label{code:javascript_composition_example}
\end{code}

Esimerkissä näytetyt funktiot ovat hyvin yksinkertaisia. Kuitenkin ne ovat täysin uudelleenkäytettäviä. Jos totuttautuu käyttämään joitain funktioita kaikissa projekteissa, voi huomata koodin kirjoittamisen tehokkuuden ja johdonmukaisuuden kasvavan.

\section{Tyyppiteoria ja joukko-oppi}
Tämä osa pyrkii kirjoittamaan matemaattisista käsitteistä käytännönläheisesti ja mahdollisesti suuresti yksinkertaistaen. Ajatuksena on, että ylemmän tason ajattelumalli välittyisi tekstissä tarkkojen määritelmien sijasta.

Joukko-oppi on matematiikan haara, joka tutkii joukkojen ominaisuuksia ja niiden välisiä suhteita. Ohjelmoinnissa joukko-oppi ilmenee yksinkertaisimmin joukko-tietorakenteiden kautta, jotka mahdollistavat uniikkien alkioiden käsittelyn tehokkaasti ja ilmaisuvoimaisesti \cite{mdn_set,mdn_set_methods}. On myös luotu ohjelmointikieli, joka perustuu kokonaan joukko-oppiin \cite{SETL_SET_LANGUAGE}.

Tyyppiteoria puolestaan on matematiikan haara, joka tutkii arvojen luokittelua, ja miten niitä käsitellään ohjelmoinnissa ja logiikassa \cite{type_theory,algebraic_data_types}.

Tyyppien ja joukkojen semanttinen ero on nyanssinen \cite{type_vs_set}. Ohjelmoinnissa puhutaan usein tyypeistä. Insinööriopintojen matematiikassa joukot ovat olleet vahvasti läsnä. Tässä osassa pääajatuksena on miettiä, kuinka asiat muutetaan funktioilla toisiksi, asiat voi mieltää kuuluvan joukkoihin tai tyyppeihin, kumpaa nyt ymmärtää paremmin. Oikea määritelmä saattaa lähestyä \glsdisp{category_theory}{kategoriateoriaa}, mutta erottaminen ei liene olevan tällä tasolla relevanttia.

Ihmisillä on kuitenkin kyky hahmottaa asioita joukkoina (tai tyyppeinä). Koko maailman voi mallintaa niillä (Kuva \ref{fig:fruit_set}) \cite{algebraic_data_types}.

\begin{figure}[ht]
    \centering
    \AltText{Omena, päärynä ja banaani.}{\includegraphics[width=10cm]{HEDELMIA}}
    \caption{Mihin joukkoon tai tyyppiin nämä saattaisivat kuulua \cite{viljofruits}?}
    \label{fig:fruit_set}
\end{figure}

Koska koko maailman voi mallintaa joukkoina (tai tyyppeinä), voidaan niillä mallintaa myös ohjelmia.


\subsection{Joukko-tietorakenne}

Ohjelmoinnissa joukko-tietorakennetta käytetään hyvin laajasti, ja sen merkitys on keskeinen monissa ohjelmointikielissä, kuten Pythonissa, Javassa ja JavaScriptissä. Joukko-tietorakenne mahdollistaa alkioiden tallentamisen ilman duplikaatteja. Sen operaatiot, kuten unioni, leikkaus ja erotus, ovat suoraan peräisin matemaattisesta joukko-opista. Joukko-tietorakenteen tehokkuus ja yksinkertaisuus tekevät siitä tehokkaan työkalun monenlaisissa sovelluksissa. \citep{mdn_set,ecma_spec}

Joukko-tietorakenne kuvaa joukko-opin joukkoa. Joukko-tietorakenteessa säilytetään kokoelma uniikkeja alkioita. Joukko-tietorakenne on ohjelmoinnissa mieluinen työkalu sen tehokkuuden ansiosta. \citep{ecma_spec}

Kokemusten mukaan ohjelmoija ei usein mieti joukko-tietorakennetta matemaattisista lähtökohdista. Tästä kuitenkin voisi olla hyötyä, sillä Joukko-opin periaatteiden ymmärtäminen ja soveltaminen ohjelmoinnissa ei ainoastaan tehostuskeino, vaan myös laajentaa näkemystä siitä, miten abstraktit matemaattiset konseptit voidaan muuttaa konkreettisiksi työkaluiksi. Tämä yhdistelmä tarjoaa ohjelmoijille sekä teoreettisen että käytännöllisen perustan, joka on arvokas monissa erilaisissa sovelluksissa ja ongelmanratkaisutilanteissa \cite{bartosz_category_for_progamers}.

\subsection{Funktioilla joukoista toisiin}
Kun ohjelmoinnissa pohditaan joukko-opin hyödyntämistä, kehittyy \glsdisp{language_agnostic}{kieliriippumaton} kyky kirjoittaa koodia. Kun miettii, mitä ylipäätään voi laskea, pääsee ohjelmoinnissa omalle abstraktille tasolle. \citep{Tan2004,BlellochHarper2015}

Matemaattisesti joukkojen välille voidaan tehdä kuvauksia, tai toisin sanoen funktioita. Merkintä funktiosta, joka \textquote{muuttaa} joukon $A$ joukoksi $B$, on esimerkiksi yksinkertaisesti $A \to B$ \cite{mellin2005joukkooppi}. Funktio muuttaa minkä tahansa joukon $A$ alkion joksikin joukon $B$ alkioksi.

Funktionaalisessa ohjelmoinnissa puhtaiden funktioiden voidaan ajatella aina olevan rinnastettaessa jonkin joukon muutokseksi johonkin toiseen joukkoon (Esimerkkikoodi \ref{code:ts_set_example}) \cite{bartosz_category_for_progamers}. (Toisaalta funktioiden argumentit ja palautusarvot ovat aina joukkoja riippumatta siitä onko funktio puhdas vai ei. Ei-puhtaat funktiot eivät kuitenkaan välttämättä luo yksikäsitteisiä muutoksia joukkojen alkioille.)

\begin{code}
    \begin{minted}{typescript}
type SetA = 'a' | 'b' | 'c'
type SetB = 1 | 2 | 3
const aToB = (a: SetA): SetB => {
    switch (a) {
        case 'a': return 1
        case 'b': return 2
        case 'c': return 3
    }
}  
\end{minted}
    \caption{Havainnollistava funktio, joka muuttaa joukon A, \{a,b,c\}, joukoksi B, \{1, 2, 3\}}
    \label{code:ts_set_example}
\end{code}

Tarkastellaan konkreettisempaa esimerkkiä lähtien koodiesimerkistä \ref{code:ts_set_theory_3}.

\begin{code}
    \begin{minted}{typescript}
const divide = (x:number) => (y:number) => x / y
\end{minted}
    \caption{Funktio, joka jakaa numeron toisella}
    \label{code:ts_set_theory_3}
\end{code}



Funktio ottaa kaksi numeroa, ja jakaa ensimmäisen toisella. Argumentit on tyypitetty JavaScript-numeroiksi. Jos kyseistä funktiota haluaa käyttää, ohjelmoijan tulisi tietää, että nollalla ei voi jakaa, ja että jos niin kuitenkin tehdään, JavaScript palauttaa arvon \mintinline{javascript}|Infinity| (tai erityisemmässä tapauksessa \mintinline{javascript}|NaN|) (Koodiesimerkki \ref{code:ts_division}).


\begin{code}
    \begin{minted}{typescript}
console.log(1/0)  // Infinity
console.log(1/-0) // -Infinity
console.log(0/0)  // NaN
\end{minted}
    \caption{Jakamisoperaattorin toiminta JavaScriptissä}
    \label{code:ts_division}
\end{code}

Jos jakolaskufunktion kehitykseen halutaan soveltaa joukko-oppia, voi olla mieluista määritellä argumenteille spesifimmät joukot, sillä joukkona JavaScriptin \mintinline{javascript}|number| sisältää epämieluisia arvoja laskujen kannalta (\mintinline{javascript}|[NaN, Infinity, -Infinity]|).

Voidaan luoda jakolaskua varten joukot \mintinline{javascript}|Rational| ja \mintinline{javascript}|NonZero|, ja ohjelmoida ne TypeScript tyypeiksi (Koodiesimerkki \ref{code:ts_type_mayhem}).

\begin{code}
    \begin{minted}{typescript}
const RationalSymbol = Symbol('Rational')
const NonZeroSymbol = Symbol('NonZero')

type Rational = number & { [RationalSymbol]: "brand" }
type NonZero = Rational & { [NonZeroSymbol]: "brand" }

const isRational = (x: number): x is Rational => 
    ![Infinity, -Infinity, NaN].includes(x)
const isNonZero = (x: number): x is NonZero => 
    isRational(x) && x !== 0
\end{minted}
    \caption{Rational ja NonZero joukkojen määritelmät TypeScriptissä}
    \label{code:ts_type_mayhem}
\end{code}

Ajatuksena on, että joukko \mintinline{javascript}|Rational| sisältää kaikki JavaScriptin tukemat rationaaliluvut, ja että joukko \mintinline{javascript}|NonZero| muuten sama paitsi ilman nollaa (Kuva \ref{fig:set-esimerkki}).



\begin{figure}[ht]
    \centering
    \AltText{Joukko NonZero on joukon Rational aito osajoukko. Joukko Rational on joukon number aito osajoukko.}{\includegraphics[width=7.1cm]{set_esimerkki}}
    \caption{Joukot NonZero, Rational ja number esitettynä osajoukkoina}
    \label{fig:set-esimerkki}
\end{figure}


Jos jakamisesimerkkiin määritetään joukot funktion argumenttien oikeiksi rajoitteiksi, ei funktiota voi käyttää väärin. TypeScript ei anna ajaa koodia, jos arvojen ei ole varmistettu kuuluvan oikeisiin joukkoihin (Koodiesimerkki \ref{code:ts_set_theory_6}) \cite{typsecript_website}.

\begin{code}
    \begin{minted}{typescript}
const divide = (x: Rational) => (y: NonZero) => x / y as Rational
\end{minted}
    \caption{Korrekti versio. Funktiolle annettavat arvot on varmistettava kuuluvan oikeisiin joukkoihin jossain vaiheessa ennen funktioon syöttöä käyttäen määriteltyjä \mintinline{javascript}|isRational| ja \mintinline{javascript}|isNonZero| funktioita }
    \label{code:ts_set_theory_6}
\end{code}

Aiempaan esimerkkiin voisi päätyä vaikkei olisi koskaan kuullut joukko-opista. Päällimmäisenä ajatuksena on kuitenkin se, että jos syötteitä ja palautusarvoja miettii joukkojen kannalta, syntyy luotettavampaa ohjelmakoodia ja vähemmän yllätyksiä. Tyypit ovat jo itsessään joukkoja, mutta niiden yksityiskohdat riippuvat käytössä olevasta, esimerkiksi ohjelmointikielen, tyyppijärjestelmästä \cite{type_theory}. Siksi ajatuksen tasolla joukkojen ajattelu on hyödyllistä.

Argumentit ja palautusarvot (tyypit) voi tietysti mieltää joukkoina myös \glsdisp{composed_function}{yhdistetyissä funktioissa} (Kuva \ref{fig:function_composition_in_sets}). Vaikka ohjelmoinnissa joukko-oppia ei rinnasteta ainoastaan funktionaaliseen ohjelmointiin,  joukko-opin ajattelu on siinä vähintäänkin tervetullutta.

\begin{figure}[htbp]
    \centering
    \begin{tikzpicture}
        % Ellipses for sets A, B, and C with background colors
        \draw[thick, ellipse, fill=green!20] (0,0.875) ellipse (0.8cm and 1.75cm);  % Set A with light blue background
        \draw[thick, ellipse, fill=blue!20] (3,0.875) ellipse (0.8cm and 1.75cm); % Set B with light green background
        \draw[thick, ellipse, fill=red!20] (6,0.875) ellipse (0.8cm and 1.75cm);   % Set C with light red background

        % Nodes for set A (without borders)
        \node[draw=none] (A1) at (0,2) {1};
        \node[draw=none] (A2) at (0,1.25) {2};
        \node[draw=none] (A3) at (0,0.5) {3};
        \node[draw=none] (A4) at (0,-0.25) {4};

        % Nodes for set B (without borders)
        \node[draw=none] (B1) at (3,2) {A};
        \node[draw=none] (B2) at (3,1.25) {B};
        \node[draw=none] (B3) at (3,0.5) {C};
        \node[draw=none] (B4) at (3,-0.25) {D};

        % Nodes for set C (without borders)
        \node[draw=none] (C1) at (6,2) {$\alpha$};
        \node[draw=none] (C2) at (6,1.25) {$\beta$};
        \node[draw=none] (C3) at (6,0.5) {$\gamma$};
        \node[draw=none] (C4) at (6,-0.25) {$\delta$};

        % Arrows from A to B
        \draw[->] (A1) -- (B1);
        \draw[->] (A2) -- (B3);
        \draw[->] (A3) -- (B2);
        \draw[->] (A4) -- (B4);

        % Arrows from B to C
        \draw[->] (B1) -- (C2);
        \draw[->] (B2) -- (C1);
        \draw[->] (B3) -- (C3);
        \draw[->] (B4) -- (C4);

        % Labels for sets
        \node[draw=none, text=green!50!black] at (-.5,3.25) {$A = \{1,2,3,4\}$};
        \node[draw=none, text=blue] at (3,3.25) {$B = \{A,B,C,D\}$};
        \node[draw=none, text=red] at (6.5,3.25) {$C = \{\alpha, \beta, \gamma, \delta\}$};

        % Arrow to new diagram
        \draw[->, line width=1.5mm, black] (3,-1.25) -- (3,-2.25) node[midway, right] {};

        % Corrected ellipses for sets A and C below the original picture
        \draw[thick, ellipse, fill=green!20] (0,-4.125) ellipse (0.8cm and 1.75cm);  % Set A with light blue background
        \draw[thick, ellipse, fill=red!20] (6,-4.125) ellipse (0.8cm and 1.75cm);   % Set C with light red background

        % Corrected nodes for set A (without borders) - Lower position
        \node[draw=none] (A1b) at (0,-3) {1};
        \node[draw=none] (A2b) at (0,-3.75) {2};
        \node[draw=none] (A3b) at (0,-4.5) {3};
        \node[draw=none] (A4b) at (0,-5.25) {4};

        % Corrected nodes for set C (without borders) - Lower position
        \node[draw=none] (C1b) at (6,-3) {$\alpha$};
        \node[draw=none] (C2b) at (6,-3.75) {$\beta$};
        \node[draw=none] (C3b) at (6,-4.5) {$\gamma$};
        \node[draw=none] (C4b) at (6,-5.25) {$\delta$};

        % Direct arrows from A to C (composition)
        \draw[->] (A1b) -- (C2b);
        \draw[->] (A2b) -- (C3b);
        \draw[->] (A3b) -- (C1b);
        \draw[->] (A4b) -- (C4b);
    \end{tikzpicture}
    \vspace{10pt}
    \caption{Funktioiden yhdistäminen joukko-opin näkökulmasta. Funktiot A$\,\to\,$B ja B$\,\to\,$C yhdistämällä saadaan A$\,\to\,$C. Älykkäällä kääntäjällä voi poistaa turhia laskutoimituksia. Keskimmäisen joukon evaluointi on käytännössä turhaa, jos funktiot ovat puhtaita}
    \label{fig:function_composition_in_sets}
\end{figure}


\subsection{Kategoriateoria}


Joukko-oppi ei ole ainoa ohjelmointiin sovellettava matemaattinen osa-alue. Funktionaaliselle ohjelmoinnille oleellisin ja syvällisin osa-alue on \gls{category_theory} \cite{bartosz_category_for_progamers,promises-spec-94,dear_functional_bros}.


Kategoriateoria tutkii matemaattisia rakenteita ja niiden välisiä suhteita erittäin abstraktilla tasolla. Se keskittyy objektien (kuten joukkojen tai tyyppien) ja niiden välillä olevien morfismien (funktioiden tai muunnosten) tutkimiseen \cite{bartosz_category_for_progamers}. Kategoriateorian avulla voidaan ymmärtää ja formalisoida monimutkaisia matemaattisia ja ohjelmallisia rakenteita. \cite{bartosz_category_for_progamers,promises-spec-94,category_theory}
Kategoriateoriaa ohjelmoijille opettavan Bartosz Milewskin kokemusten mukaan, vaikka algebra tai kalkyyli olisi ohjelmoijalle ajatuksena hirveää, kategoriateoria on kuitenkin erityisen mieluisaa sen suuresta teoreettisuudesta huolimatta. \cite[9]{milewski2017category}.

Kategoriateoriassa peruspalikat ovat kategoriat \cite[9]{milewski2017category}. Kategoria on rakenne, joka koostuu morfismeista. Esimerkiksi joukot ja funktiot koostavat kategorian, jossa joukot ovat kategorian objektit, ja funktiot kategorian morfismit \cite{category_theory}. Kategoriateoriaan ei perehdytä tarkemmin tässä työssä Monadi-rakennetta lukuun ottamatta. Ei ole käytännönläheistä tuoda tavattoman teoreettisia käsitteitä työhön, jossa tavoitteena on etsiä jonkin kaltaista konkretiaa.


\section{Monadi}

Monadi tulee esille funktionaalisessa ohjelmoinnissa erittäin usein. Se on alun perin teoreettinen rakenne, joka on löytynyt 1900-luvun puolivälin jälkeen kategoriateoriaa tutkivien matemaatikkojen ansiosta \cite{Beck2003}. Monadin hyöty ohjelmoinnissa ilmeni vasta myöhemmin 1990-luvulla \cite{computerphile_monad}.

Ohjelmoinnissa monadilla voidaan hallita sivuvaikutuksia, ketjuttaa operaatioita ja yhtenäistää abstraktioita \cite{computerphile_monad,bartosz_category_for_progamers_10,stackoverflow_what_monad}. Ohjelmoinnissa monadi on kieliagnostinen käsite.

Tässä osassa käydään läpi, mitä ovat yleiset monadit, miten implementoida monadi, miten niitä käytetään, ja lopulta pyritään tiivistämään miten monadeja voi ajatella ohjelmoinnissa.

\subsection{Yleiset monadit}

\texttt{Promise, Maybe, Either, Result, State, IO, ja Lista (Array)}. Nämä ovat esimerkkejä yleisesti käytetyistä monadeista. \cite{monad_wikipedia,bartosz_category_for_progamers_10}

Pintapuolisesti tarkasteltuna esitetyt käsitteet vaikuttavat hyvin erilaisilta, sillä niiden käyttötarkoitukset ja rakenteet vaihtelevat merkittävästi. Esimerkiksi \texttt{Promise} käsittelee asynkronisia operaatioita, kun taas \texttt{Maybe} kuvastaa mahdollisia arvoja, jotka voivat olla olemassa tai olla olematta. \texttt{Either} ja \texttt{Result} tarjoavat mekanismin virheiden käsittelyyn, ja \texttt{State} mahdollistaa tila-arvojen kantamisen funktioiden läpi. \texttt{IO} käsittelee sivuvaikutuksia, ja \texttt{Lista (Array)} tarjoaa toistuvien arvojen käsittelyn.

Kuitenkin näitä kaikkia yhdistää yhteinen abstraktio. Ne kaikki toteuttavat tietyt perustoiminnot \texttt{bind} (\texttt{flatMap/chain}) ja \texttt{of} (\texttt{unit/pure}). Näiden operaatioiden avulla monadit voivat kapseloida laskennan tilan tai sivuvaikutuksen, tarjoten yhdenmukaisen tavan yhdistää operaatioita ja siirtää tietoa laskentaketjujen läpi.

JavaScriptissä \mintinline{js}|Array| toteuttaa toiminnon \texttt{bind} \mintinline{js}|Array.prototype.flatMap|-metodilla, ja \texttt{of} toiminnon staattisella \mintinline{js}|Array.of|-metodilla \cite{stackoverflow_flatmap_monad,stackoverflow_js_array_monad}.


JavaScriptin \mintinline{js}|Promise| toteuttaa toiminnon \texttt{bind} \mintinline{javascript}|Promise.prototype.then|-metodilla, ja \texttt{of} toiminnon staattisella \mintinline{javascript}|Promise.resolve| ja \mintinline{javascript}|Promise.reject| metodeilla \cite{read-it-later-11481,stackoverflow:why_monad,promises-spec-94}.

Vaikka JavaScriptin \mintinline{js}|Promise| ja \mintinline{js}|Array| toteuttavatkin monadin \texttt{bind} ja \texttt{of} toiminnot, se ei tee niistä automaattisesti monadeja kategoriateorian tai funktionaalisen ohjelmoinnin sääntöjen mukaan \cite{promises-spec-94,stackoverflow:why_monad,stackoverflow_js_array_monad}. Kehittäjä Jake Archibaldin mukaan akateemisten käsitteiden seuraaminen puhtaasti ei kuitenkaan usein johda käytännölisiin ratkaisuihin \cite{pennane_fp_gist}. Tästä syystä on kohtuullista puhua näistä rakenteista monadeina, vaikkeivat ne sitä puhtaasti olisikaan.


\subsection{Monadin implementointi}

Konkreettisesti monadilla tulee siis olla kaksi toimintoa: \texttt{of} ja \texttt{bind}. Toiminto \texttt{of} käärii arvoja monadiin. Toiminnon voidaan ajatella olevan monadin konstruktori. \citep{stackoverflow_what_monad}

\texttt{Bind}-toiminnolla muutetaan monadiin käärittyä arvoa antamalla toiminnolle funktio, joka palauttaa toisen saman tyypin monadin. \cite{stackoverflow_what_monad} \texttt{Bind}-toiminto on käytännössä sama kuin tunnetumpi \texttt{map}-toiminto. Erona on, että \texttt{map} palauttaa muokatun arvon, joka on kääritty samaan monadiin, kun taas \texttt{bind} mahdollistaa monadin sisällä olevan arvon muuttamisen ja sen ketjuttamisen muihin monadisiin operaatioihin.

Näillä kahdella toiminnolla voidaan luoda rakenteita, jotka käyttäytyvät kuin monadi. Monadin täydet hyödyt eivät kuitenkaan ilmene, jos rakenne ei myös täytä monadien kolmea sääntöä: vasen identiteetti (left identity), oikea identiteetti (right identity), ja liitännäisyys (associativity) \cite{haskellmonadlaws}. Jos kaikki kolme sääntöä toteutuu, voidaan luoda työkaluja, joissa voidaan käyttää mitä tahansa monadia, ja silloin monadi on määritelmän mukaisesti monadi.

Osassa \textquote{Promise-tietorakenne on miltei monadi} puidaan sitä, miksei JavaScriptin \texttt{Promise} ole määritelmällisesti täydellisesti monadi, ja miksei ilman ulkoisia kirjastoja JavaScriptiin voida enää tuoda määritelmän täyttäviä monadeja. Kuitenkin rakenne, joka on kuin monadi, voi myös olla erittäin hyödyllinen. Siksi ei tarkemmin käsitellä, miten monadin määritelmän voi oikeellisesti täyttää.


\subsection{Monadin käyttäminen}

Monadeja käytetään operaatioiden ketjuttamiseen. Monadeilla voi käytännössä tehdä yhdistettyjä funktioita, joissa on mukautettu tapa hallita mitä yhdisteen funktiokutsujen välissä tapahtuu.

Käydään esimerkkinä JavaScriptin \texttt{Promise}-tietorakennetta (joka toimii kuin monadi) (Koodiesimerkki \ref{code:promise_monad_concise}).

\begin{code}
    \begin{minted}{typescript}
fetch("https://some.imaginary.api.com/v1") 
    .then(result => result.json()) 
    .then(data => data.users)
    .then(users => users.map(user => user.id))
    .catch(error => Array.of()) 
    \end{minted}
    \caption{Promise-tietorakenteella bind-operaatioden ketjutus, jossa haetaan käyttäjätietoja kuvitteellisesta ulkoisesta rajapinnasta. Jos mikään yksittäinen askel päätyy virhetilaan, palauttaa ohjelma tyhjän listan kutsumalla \mintinline{js}|Array.of|-metodia}
    \label{code:promise_monad_concise}
\end{code}

Koodiesimerkissä nähdään operaatioiden monadinen ketjutus. Jokainen \mintinline{js}|Promise.then|-kutsu palauttaa aina uuden Promise-monadin. Jos jokin yksittäinen Promise päätyy virhetilaan, ei enää seuraavia then-kutusja ajeta, vaan ketju päätyy automaattisesti \texttt{catch}-osioon.

Käytännössä koko ohjelma on siis vain yhdistetty funktio, jossa ollaan määritelty räätälöityä logiikkaa siihen, miten funktiokutsut käsitellään. Kuitenkaan \mintinline{js}|Promise.catch| ei liity suoraan monadisuuteen, eikä käytännössä oikeellisesti implementoidun monadin tarvitsisi käyttää sellaista \cite{promises-spec-94}.

Raja on häilyvä. Onko JavaScriptin Promise monadi vai ei? Matemaatikkojen mielestä monadi on \textquote{yksinkertaisesti} monoidi endofunktoreiden kategoriassa \cite{bartosz_category_for_progamers_10,monad_wikipedia,stackoverflow_what_monad}. Toisten mielestä monadi on asia, jolla on \texttt{flatMap}-metodi \cite{stackoverflow_flatmap_monad}.

Pragmaattisuuden nimissä monadi saa tässä työssä olla rakenne, jolla voidaan ketjuttaa operaatioita itsemäärittelemällä tavalla, joka mahdollisesti eroaa tavallisesta funktioiden yhdistämisestä.


\section{Pragmaattisuus funktionaalisessa ohjelmoinnissa}

Funktionaalisen ohjelmoinnin iskevin teoria on pitkälti käytynä, ja on ehkä selvää, että funktionaalinen ohjelmointi on pohjimmiltaan erittäin teoreettista. Voi pohtia, miten funktionaalinen ohjelmoinnin voi tuoda ohjelmistokehittäjän arsenaaliin käytännönläheisesti.

Kun funktionaalisen ohjelmoinnin kauneus on matemaattisissa kulmakivissä, joilla ei ole tapaa taipua, miten sen voi yhdistää säännöttömään JavaScriptiin, tai muuhun vastaavanlaiseen ohjelmointikieleen?

Mikä on se lähestymistapa, jota tarvitaan käytännönläheiseen, pragmaattiseen, funktionaaliseen ohjelmointiin?

Robert C. Martinin mukaan pragmaattinen funktionaalinen ohjelmointi on sitä, että käytetään esimerkiksi \glsdisp{immutable_data}{muuttumattomuutta} ja puhtaita funktioita silloin, kun ne parantavat ohjelman luettavuutta, ylläpidettävyyttä ja suorituskykyä, ja kuitenkin niin, ettei vaadita täydellistä sitoutumista funktionaalisen ohjelmoinnin paradigmaan. \citep{martin2017pragmaticfp}

Myös Cantarella mainitsee opinnäytetyönsä loppupuolella ajatuksen, että mitä enemmän paradigmoihin liittyviin artikkeleihin perehtyy, sitä myötämielisemmäksi tulee ajatukselle, että paradigmoja tulisi yhdistellä vapaammin. \citep[45]{cantarella_fp_haitat}

Myös kokemusten perusteella työpaikalla on havaittu funktionaalisen ohjelmoinnin toimivan kohtuullisen hyvin ripoteltuna proseduraalisen ohjelmakoodin.

Pragmaattista funktionaalista ohjelmointia vaikuttaa siis olevan se, että otetaan vain mitä tarvitaan \cite{dear_functional_bros,martin2017pragmaticfp,cantarella_fp_haitat}. Toisaalta se, mitä tarvitaan ei ole välittömästi selvää. Toisille pragmaattista on jättää kaikki matemaattinen teoria pois ohjelmoinnista heti kättelyssä, toisille kategoriateorian rakenteiden käyttäminen, kun ne oikeasti sopivat mallintamaan ongelmaa \cite{holvikari2021category,martin2017pragmaticfp}.

Tässä insinöörityössä tarkempia pragmaattisia funktionaalisen ohjelmoinnin käytäntöjä pyritään tuomaan esille neljännessä osassa \textquote{Funktionaalisia käytänteitä ei-funktionaalisissa ohjelmointikielissä}, ja siitä eteenpäin.

\section{Funktionaalinen ohjelmointi \& TypeScript}
JavaScript ei ole funktionaalinen ohjelmointikieli, vaikka se funktionaalisuutta tukeekin \cite{is_reduce_bad}. TypeScriptin tyyppijärjestelmällä funktionaalisesta ohjelmoinnista saa jo mieluisampaa \cite{holvikari2021category}.

JavaScriptia halutaan käyttää kaikkeen ohjelmointiin, vaikka usein erinäköisiin tehtäviin olisi parempiakin ohjelmointikieliä tarjolla. Voi itse päättää, onko pragmaattista käyttää kaikkeen samaa työkalua, ja säästää opiskelussa, vai käyttää aikaa uuden sopivamman työkalun löytämiseen, ja säästää toteuttamisessa.


\subsection{Kielen sopivuus funktionaaliseen ohjelmointiin}

Vaikka TypeScript ei ole pohjimmiltaan funktionaalinen ohjelmointikieli, on sillä ominaisuuksia, joilla funktionaalisesta ohjelmoinnista tulee verraten helppoa.

TypeScriptissä on käytännössä kolme eri tapaa määrittää funktioita (Esimerkkikoodi \ref{code:javascript_function_types}). Nuolifunktioilla voidaan harjoittaa funktionaaliselle ohjelmoinnille tyypillistä argumenttien yksi kerrallaan antamista.


\begin{code}
    \begin{minted}{typescript}
function rnd(min: number, max: number) { 
	return Math.floor(Math.random() * max) + min;
}
const rnd2 = (min, max) => Math.floor(Math.random() * max) + min;
const rnd3 = min => max => Math.floor(Math.random() * max) + min;
\end{minted}
    \caption{Kolme eri tapaa kirjoittaa funktio JavaScriptissä \cite{okhravi-g-discussion}. Funktiomäärittely, funktioilmaus ja osittain sovellettava funktioilmaus}
    \label{code:javascript_function_types}
\end{code}

Koska funktioita voi käyttää JavaScriptissä kuin muuttujia (first-class citizen), voi niitä antaa argumentteina toisille funktioille (Koodiesimerkki \ref{code:js_first_class}). Tämä mahdollistaa funktionaaliselle ohjelmoinnille tyypillisten \glsdisp{higher_order_function}{korkeamman asteen funktioiden} luomisen, ja käyttämisen.

\begin{code}
    \begin{minted}{javascript}
function A()  { return "moi"; }
function B(f) { return f(); }
B(A) // "moi"
\end{minted}
    \caption{Funktioiden ensiluokkaisuus JavaScriptissä. Funktiolle B voidaan antaa funktio argumenttina}
    \label{code:js_first_class}
\end{code}

Funktionaalisen ohjelmoinnin piireissä TypeScriptin tyyppijärjestelmää arvostetaan myös, koska sillä voidaan rakentaa algebrallisia datatyyppejä, jotka ovat yleisesti puhtaissa funktionaalisissa ohjelmointikielissä läsnä ja suuressa roolissa \cite{holvikari2021category}.


\subsection{Algebralliset datatyypit}

TypeScriptin tyypit ovat mieluisia funktionaaliselle ohjelmoinnille. TypeScript mahdollistaa funktionaaliselle ohjelmoinnille olennaisten algebrallisten datatyyppien määrittelemistä \cite{holvikari2021category}.

Algebralliset datatyypit ovat osa funktionaalisen ohjelmoinnin maailmaa. Niiden käyttämisen on huomattu olevan järkevää ja luontevaa \cite{holvikari2021category,hickey_maybe_not,algebraic_data_types}

Yksinkertaisesti algebralliset datatyypit ovat keino määritellä ja käsitellä monimutkaisia datastruktuureja, jotka koostuvat perustietotyypeistä yhdistämällä niitä eri tavoin. Tyypillisesti algebralliset datatyypit jaetaan kahteen päätyyppiin: tulo- ja summatyypit.

\textbf{Tulotyyppi} (product type) kuvaa tietorakenteita, joissa yhdistetään useita arvoja yhdeksi kokonaisuudeksi, esimerkiksi luokkien tai tietueiden (records) avulla.
TypeScriptissä tulotyyppejä voi määritellä käyttäen objekteja tai Record-tyyppejä (Koodiesimerkki \ref{code:ts_product_type}). \citep{algebraic_data_types,holvikari2021category}

\begin{code}
    \begin{minted}{typescript}
// Tyypin mahdollisten erilaisten alkioiden määrä: värimahdollisuudet * istuinmahdollisuudet -> tulotyyppi
type Car = {type: "car", color: string; seats: number;}
\end{minted}
    \caption{Tulotyyppi-esimerkki TypeScriptissä. Tulotyyppi tulee siitä, että tyypin osien määrän voi kertoa keskenään saadakseen tyypin kokonaisen permutaatioden määrän}
    \label{code:ts_product_type}
\end{code}

\textbf{Summatyyppi} (sum type) mahdollistaa arvon, joka voi olla yksi useista ennalta määritellyistä tyypeistä.TypeScriptissä summatyyppejä voi määritellä käyttäen diskriminoituja unioneja (discriminated union) (Koodiesimerkki \ref{code:ts_sum_type}). \citep{algebraic_data_types,holvikari2021category}

\begin{code}
    \begin{minted}{typescript}
type Car = {type: "car", color: string; seats: number;}
type MotorBike = {type: "motorbike", top_speed: number; wheelieability:number;}
// Tyypin mahdollisten erilaisten alkioiden määrä: autojen määrä + moottoripyörien määrä -> summatyyppi
type Vehicle = Car | MotorBike
\end{minted}
    \caption{Summatyyppi-esimerkki TypeScriptissä. Summatyyppi tulee siitä, että tyypin osien määrän voi summata keskenään saadakseen tyypin kokonaisen permutaatioden määrän}
    \label{code:ts_sum_type}
\end{code}


Vaikka TypeScriptin tyyppijärjestelmä on vankkarakenteinen, ei se varmista kaikista monimutkaisimpien tyyppien \glsdisp{correctness}{oikeellisuutta} \cite{holvikari2021category}. Sellaisiin tarvitaan jokin pohjimmiltaan funktionaalinen ohjelmointikieli, kuten Haskell \cite{holvikari2021category}.

\subsection{Kirjastot}

Funktionaalista ohjelmointia on pyritty pitämään elossa erinäköisin kirjastoin, joissa hyväksikäytetään kielen ominaisuuksia, ja luodaan kirjastoja, joiden käyttäminen vaatii ohjelmoijalta täysin uuden syntaksin opettelemista JavaScript syntaksin lisäksi \cite{ramda,sanctuary,crocks,fpts,holvikari2021category}.

Kokemusten perusteella syntaksierot ovat todella hidastava tekijä. Ohjelmoijat olettavat ohjelmoivansa TypeScript-projektissa TypeScriptiä. Tai yleisemmin sanottuna: ohjelmoijat olettavat ohjelmoivansa sitä ohjelmointikieltä mitä ohjelmoivat.

Funktionaalisen ohjelmoinnin kirjastoilla halutaan tuoda \gls{ts}-ekosysteemiin sitä, mitä siitä puuttuu. Kuitenkin on kirjasto kuinka hyvälaatuinen tahansa, tulee TypeScriptin tyyppijärjestelmän rajoitteet kuitenkin lopulta vastaan \cite{holvikari2021category}.

\subsection{Promise-tietorakenne on miltei monadi}

JavaScriptissä oleva Promise-tietorakenne toimii miltei kuin monadi, muttei ole sitä määritelmällisesti matematiikan kategoriateorian mukaan \cite{stackoverflow:why_monad,promises-spec-94}. Tässä osassa puidaan miksi se on merkittävää, miksi se pääsi tapahtumaan, ja mitä merkitystä sillä edes on, että JavaScriptissä ei ole standardoitua tapaa luoda monadeja.

JavaScriptin kehitys perustuu vahvasti takaperin yhteensopivuuteen, mikä estää merkittäviä muutoksia kielen rakenteisiin ilman riskiä rikkoa olemassa olevia verkkosivuja \cite{prototype_library_trends}. Selaimet eivät voi myöskään tehdä omin päin rikkovia muutoksia \cite{against_self_closing_tags,proposal-joint-iteration}. Jos vierailee lääkärin vastaanoton verkkosivulla selvittääkseen aukioloajat, kumpi selain on paras: se, joka näyttää vastaanoton aukioloajat, vai se, joka ei näytä mitään, koska on ottanut uuden JavaScript-version käyttöön?

Kuitenkin JavaScript kehittyy. Vuonna 2015 Promise-tietorakenne lisättiin JavaScriptiin ratkaisemaan callback-funktioiden hallinnan ongelmia asynkronisessa ohjelmoinnissa \cite{mdn_promise,callbackhell}. Promise-tietorakenteen juuret ovat funktionaalisessa ohjelmoinnissa. Promise-tietorakenne on rakenteeltaan sellainen, että sen voi kuvata monadina \cite{promises-spec-94,stackoverflow:why_monad}.

Kun Promise-tietorakennetta oltiin tuomassa osaksi JavaScriptia, keskusteltiin tulisiko JavaScriptin Promise-tietorakenteen noudattaa \glsdisp{monad}{monadirakennetta} \cite{promises-spec-94}. Keskustelu jäi tuloksettomaksi, vaikka muutokset implementointiehdotukseen olisivat olleet pienet monadin sääntöjen pitämiseksi.

Yksityiskohdat eivät liene merkittäviä, mutta lopputulos on se, ettei JavaScriptin Promise-tietorakennetta saatu täyttämään kaikkia monadien sääntöjä.

Mahdollisuus tuoda monadit osaksi JavaScript standardia on todennäköisesti menetetty \cite{proposal-joint-iteration,prototype_library_trends}. Ohjelmoijat joutuvat itse valitsemaan, miten luoda ja käyttää monadeja JavaScriptissä. Jos sekä Array että Promise olisivat yhtenäisesti monadeja, niiden yhteiskäyttö olisi suoraviivaisempaa, ja uusien monadien implementointi olisi johdonmukaisempaa suoraan JavaScriptissä kielen tasolla. Nykytilanne pakottaa ohjelmoijat improvisoimaan.


% Material and Methods
%\clearpage%if the chapter heading starts close to bottom of the page, force a line break and remove the vertical vspace
\vspace{21.5pt}
\chapter{Funktionaalisia käytänteitä ei-funktionaalisissa ohjelmointikielissä}

Tässä osassa pyritään esittämään kaikista oleellisin, mitä funktionaalisella ohjelmoinnilla voidaan saada aikaan. Teoreettisen perustan ja kokemusten avulla tuodaan esille mitä funktionaalinen ohjelmointi on koetusti parhaimmillaan.

Ajatuksena on, että kun ohjelma on mallinnettu oikein, ei ole suurta merkitystä ohjelmoidaanko olio-ohjelmointia, imperatiivista ohjelmointia, vai funktionaalista ohjelmointia. Olio-ohjelmoinnissa mallinnus saatettaisiin hoitaa rajapinnoilla (interface). Funktionaalisen ohjelmoinnin tapa mallintaa on käyttää tyyppejä ja abstrakteja (algebrallisia) tietorakenteita.

Tuodaan ilmi millaisia käytänteitä voidaan ripotella ohjelmointiin, vaikka ohjelmointikieli ei varsinaisesti olisi funktionaalinen.

Pyritään herättämään ajatuksia milloin ja miten funktioita kannattaa käyttää. Esitetään kuinka suuressa roolissa ongelmien mallintaminen on. Kun ongelmien mallit pureskellaan oikeisiin joukkoihin, saadaan turvallisempaa ohjelmakoodia. Silloin on myös helpompi kirjoittaa deklaratiivista koodia, joka on funktionaaliselle ohjelmoinnin ominaista.



\section{Opi kerran, käytä kaikkialla}

Funktionaalinen ohjelmointi kannustaa asioiden uudelleenkäytettävyyteen. Funktionaalisessa ohjelmoinnissa etsitään yleisiä teemoja, ja viedään toistuvia malleja funktioiksi, jotta voidaan kirjoittaa funktioiden nimiä syntaksin sijasta.


Esimerkiksi olio-ohjelmoinnissa, tai muuten imperatiivisessa ohjelmoinnissa, for-silmukalla toteutettavat algoritmit voidaan yleistää funktioiksi (Kuva \ref{fig:ramdacmds}). Onko helpompaa opetella ääretön määrä funktioita, vai yksi kielirakenne: for-silmukka?

\begin{figure}[ht]
    \centering
    \[
        \begin{array}{rl}
            \left.
            \begin{array}{l}
                filter, without, find, findIndex, findLast, findLastIndex \\
                map, mapIndexed, mapIndexedRight, chain, concat, zip      \\
                reduce, reduceRight, scan, partition, uniq, forEach       \\
                slice, drop, take, dropWhile, takeWhile, zip, zipWith     \\
                all, any, none
            \end{array}
            \right] \quad \text{for loop}
            \\
        \end{array}
    \]
    \caption{Ramda.js-kirjaston funktioita listojen käsittelyyn ohjelmoinnissa \cite{ramda:docs}. Kaikki nämä voitaisiin korvata for-silmukoilla imperatiivisessa paradigmassa}
    \label{fig:ramdacmds}
\end{figure}

Nimet ovat syntaksiriippumattomia. Samoja nimiä voi käyttää ohjelmointikielestä riippumatta. Nimillä voi kuvata oikean maailman asioita. Siksi voi perustella, että kun tietyn logiikan aina laittaa tietyn nimen taakse funktioksi, ymmärtää sen aina vaikka ohjelmointikieli vaihtuisi. Jos muistaa idean, ei implementaation yksityiskohdilla ole merkitystä.

Jos asia on tarpeeksi monimutkainen ja sen piilottaa nimen taakse, hyötyarvon määrittäminen voi olla haastavaa. Mitä monimutkaisempaa asiaa yritetään enkoodata yksittäiseen nimeen, sitä vaikeampaa se on lukijan ymmärtää, ja myös sitä vaikeampaa sille on löytää uudelleenkäyttötilanteita. Pienikin nyanssiero funktion implementaatiossa totuttuun implementaatioon voi mitätöidä opitun hyödyn eri ympäristöissä.

Yksinkertaisiakin funktioita voi olla vaikea hahmottaa koodista. Mikään ei pakota ohjelmoijaa nimeämään funktiota samalla tavalla kuin miten sama funktio on nimetty jossain toisessa projektissa tai ohjelmointikielessä. Nimeämiskäytänteiden noudattaminen ei kuitenkaan ole helppoa. Käytänteitä voi olla vaikea löytää, tai niitä voi olla useita \cite{ramda:docs}.

Kannattanee käyttää ohjelmointikielen sisäänrakennettuja funktioita, ja metodeja aina kun mahdollista. Näin voi maksimoida sen, että koodi ymmärretään. Jos kyse on jostakin projektin sisäisestä kirjastosta, voi olla hyödyllistä lisätä funktiolle kommentti, jos se tunnetaan muissa ohjelmointiympäristöissä jollain toisella nimellä.

Jos on kyse jostain toimintaympäristökohtaisesta funktiosta, eikä perustavanlaatuisesta yleispätevästä funktiosta, on tärkeää, että funktio nimetään mahdollisimman kuvaavasti. Esimerkiksi sen sijaan, että nimeää funktion nimellä \textquote{processUsers}, voi miettiä tarvitseeko funktion kutsupaikoissa ymmärtää enemmän siitä mitä funktio tekee? Jos kontekstista on selvää, mitä funktio tekee, nimen voi jättää sellaiseksi kuin on. Jos konteksti ei avaa toimintaa enemmän, \textquote{processUsers} on erittäin ympäripyöreä nimi.

\section{Yhdistetyt funktiot}

Kun funktioita on saatu kasattua työkalupakkiin, niitä voi yhdistellä erilaisten tarpeiden mukaan. Kerran kirjoitetut funktiot ovat monikäyttöisiä ja niitä voidaan ketjuttaa keskenään, jolloin monimutkaisemmat operaatiot rakentuvat valikoiduista palikoista.

Yhdistettyjen funktioiden avulla voidaan saavuttaa korkea taso ohjelmoinnin joustavuudessa. Ad hoc -funktioiden, eli tilanteeseen nopeasti mukautettujen funktioiden, kirjoittaminen on vaivatonta, ja niiden lisääminen osaksi olemassa olevia ketjuja onnistuu helposti. Tämä mahdollistaa nopean iteroinnin ja mukautumisen uusiin tarpeisiin ilman merkittäviä muutoksia alkuperäiseen koodiin.

Yhdistetyt funktiot eivät kuitenkaan automaattisesti takaa koodin luettavuutta. Liiallinen abstraktio voi tehdä koodista vaikeasti seurattavaa, erityisesti jos funktion yhdisteleminen tuo mukanaan monimutkaisia käsitteitä, jotka eivät ole suoraan yhteydessä ratkaistavaan ongelmaan. Turha abstraktio voi johtaa tilanteeseen, jossa koodin toiminnan ymmärtämiseksi täytyy perehtyä moniin välivaiheisiin ja ylimääräisiin konsepteihin, vaikka ongelma itsessään olisi yksinkertainen.

Esimerkkikoodi \ref{code:js_garbage_pipe} näyttää tilanteen, jossa Ramda.js-kirjastoa on käytetty yhdistetyn funktion rakentamiseen. Vaikka funktio toimii, sen ymmärtämiseen tarvitsee paljon ylimääräistä tietoa.

\begin{code}
    \begin{minted}{typescript}
const convertUsersToLeaderboardUsers = (users: Array<User>, anonymizeUsers: boolean) => R.pipe(
    R.pick(["userName", "score", "id"]),
    R.ifElse(
        R.always(anonymizeUsers), 
        R.map(R.omit(["userName"])), 
        R.identity
    ),
    R.map(R.assoc("type", "leaderboardUser")),
    R.sortWith([R.descend(R.prop("score"))])
)(users)
\end{minted}
    \caption{Funktio, muuttaa listan käyttäjiä sellaisiksi, että niitä voidaan käyttää tuloslistoilla. Funktio on yhdistetty monesta funktiosta käyttämällä Ramda.js-kirjaston funktioita. Funktiossa käytetään useita muita funktioita, joiden toiminnan lukijan on tiedettävä tai arvattava}
    \label{code:js_garbage_pipe}
\end{code}

Sen sijaan esimerkkikoodissa \ref{code:js_legit_pipe} käytetään vain JavaScriptin sisäänrakennettuja työkaluja. Funktio on yleisemmin helppo lukea. Vaikka luettavuus on subjektiivista, ja perustuu pitkälti ohjelmointityylin tuttavuuteen, on todennäköisempää, että ohjelmoija ymmärtää JavaScriptin omia Array-metodeja, kuin minkäkin kirjaston funktioita, vaikka kirjaston funktiot olisikin pyritty kirjoittamaan noudattaen yleisiä käytänteitä.

\begin{code}
    \begin{minted}{typescript}
const convertUsersToLeaderboardUsers = (users: Array<User>, anonymizeUsers: boolean) => users
    .map(user => ({
        type: "leaderboardUser",
        userName: anonymizeUsers 
            ? undefined 
            : user.userName,
        score: user.score
    }))
    .sort((userA, userB) => userB.score - userA.score)
    \end{minted}
    \caption{Toiminnaltaan sama funktio, mutta käytettynä on vain JavaScriptin sisäänrakennettuja palikoita. Useampi lukija ymmärtää funktion, sillä se ei käytä ulkoisia kirjastoja}
    \label{code:js_legit_pipe}
\end{code}


On tärkeää, että abstrahoinnin taso pysyy hallittuna, eikä yhdistetty funktio tuo mukanaan lisämonimutkaisuutta. Abstraktioiden on oltava perusteltuja ja niiden tulee aidosti parantaa koodin ylläpidettävyyttä ja uudelleenkäytettävyyttä, eikä vain noudattaa jotain häilyviä teoreettisia malleja ilman konkreettisia hyötyjä.


Pitää siis pitää mielessä ohjelmointikieli, jossa toimitaan. Säännöt tulee sopeuttaa ympäristöön. Vaikka funktionaalisessa ohjelmoinnissa pyritään olla mutatoimatta dataa, tästä ei missään nimessä ole pakko pitää kiinni kynsin hampain.

\section{Puhdas funktio ei ole aina paras ratkaisu}

Funktionaalisessa ohjelmoinnissa pyritään välttämään mutatointia, eli ajonaikaisten muuttujien arvojen muuttamista. Tietorakenteiden muuttamisen sijasta pyritään luomaan tietorakenteista uusia kopioita, ja jättämään vanhat ennalleen \cite{immutablejs_immutable,hickey_persistent_2009}.

Joissain ohjelmointikielissä, tai ohjelmointikielien kirjastoissa, muuttumattomien tietorakenteiden käyttäminen on saatu varsin tehokkaaksi \cite{hickey_persistent_2009,immutablejs_immutable}. Näin ei kuitenkaan ole asian laita JavaScriptissä, tai ohjelmointikielissä yleisesti. Tiedon mutatointi on nopeaa ja kopiointi hidasta \cite{turner_trauring_copying_2020}. Näin ollen, jos kirjoitetaan funktioita, jotka palauttavat aina uusia kopioita, vanhojen mutatoimisen sijasta, saadaan usein hidasta ja tehotonta ohjelmakoodia.

Jossain vaiheessa ohjelman suoritusta tietokoneen bitit kääntyvät nollista ykkösiin, tai toisin sanoen mutatoituvat \cite{is_reduce_bad}. Niin tietokoneet toimivat, minkä vuoksi mutaatiota ei saa pelätä, kunhan sen hoitaa asianmukaisesti.

Jos funktio mutatoi ainoastaan muuttujia, joita se itse luo funktion ajon ajaksi, ei ole koodin turvallisuuden kannalta toiminnallista merkitystä mitä funktion sisällä tapahtuu. Kuitenkin tehokkuusvoitot voivat olla huomattavat versioihin, joissa mutatointia ei harjoiteta kuten seuraavat esimerkit osoittavat.

Esimerkkinä funktio, joka ottaa sisään listan avain-arvo-pareja, ja luo niistä olion, josta voi hakea arvoja avainten perusteella. Näytetään kaksi (funktionaalista) tapaa toteuttaa tämä mielivaltainen funktio.
Versio, joka ei mutatoi muuttujia missään vaiheessa (Koodiesimerkki \ref{code:js_just_mutate_immutable}), ja versio, joka mutatoi vain itse luomiaan muuttujia (Koodiesimerkki \ref{code:js_just_mutate_mutable}).

Tarkkailijan näkökulmasta molemmat funktiot ovat kuitenkin puhtaita.

\begin{code}
    \begin{minted}[highlightlines={8,14}]{javascript}
const entries = Array.from({ length: 1000 }, (_, i) => [
    `key-${i}`,
    `value-${i}`
])

const objectFromEntries = (entries) =>
    entries.reduce((acc, entry) => {
        return { ...acc, [entry[0]]: entry[1] }
    }, {})

console.time('Immutable')
objectFromEntries(entries)
console.timeEnd('Immutable')
// Immutable: 112.718017578125 ms
\end{minted}
    \caption{Funktio, joka ottaa listan avain-arvo-pareja ja luo niistä olion. Olion luonnissa ei käytetä ollenkaan mutatointia Funktion suorittaminen 1000:lle avain-arvo-parille kesti n. 113 millisekuntia}
    \label{code:js_just_mutate_immutable}
\end{code}

Täysin mutatoimattomalla versiolla kesti luoda tuhannesta avaimesta olio yli 300 kertaa kauemmin, kuin versiolla, joka mutatoi käyttämiään muuttujia (113 ms ja 0.4 ms).


\begin{code}
    \begin{minted}[highlightlines={8,15}]{javascript}
const entries = Array.from({ length: 1000 }, (_, i) => [
    `key-${i}`,
    `value-${i}`
])

const objectFromEntries = (entries) =>
    entries.reduce((acc, entry) => {
        acc[entry[0]] = entry[1]
        return acc
    }, {})
          
console.time('Mutable')
objectFromEntries(entries)
console.timeEnd('Mutable')
// Mutable: 0.35693359375 ms
\end{minted}
    \caption{Sama funktio kuin aiempi. Ainoa ero, että pareja iteroidessa luotavaa oliota mutatoidaan. Funktion suorittaminen 1000:lle avain-arvo-parille kesti vain n. 0.4 millisekuntia}
    \label{code:js_just_mutate_mutable}
\end{code}

Useimmiten ohjelman tehokkuudella ei ole paljoa merkitystä, ja luotettavuus ja luettavuus on syytä pitää etusijalla. Kuitenkin tässä tilanteessa tehokkuuserot ovat niin merkittävät, että jos 1 000:n avain-arvo-parin sijasta, olisi käytetty 100 000, tai vaikka 1 000 000 avain-arvo-paria, mutatoimaton versio ei välttämättä olisi suoriutunut tehtävästä koskaan. Näitä määriä yritettiin, mutta suoritukset lopetettiin kesken, kun alkoi näyttämään, ettei tehtävä tosiaan ollut suoriutumassa koskaan.

\section{Ongelmien mallintaminen lisää turvallisuutta}

Funktionaalinen ohjelmointi mahdollistaa deklaratiivisen ja itsensä selittävän koodin kirjoittamisen \cite{ms:fp_vs_imperative}. Oikeiden nimien löytäminen funktioille tai muille ohjelmiston osille on vaikeaa, jos ongelma on mallinnettu väärin. Tässä osiossa käsitellään pragmaattista mallinnusta pala palalta kuvitteelliselle alustalle.

Käydään läpi valheellisen mallin haitat ja oikean mallin hyödyt. Staattisen tyyppijärjestelmän avulla oikea malli poistaa ohjelmavirheitä jo käännösvaiheessa ja vähentää testaustarpeita \cite{impossiblebetter}.

Oikeellisen mallin avulla koodi on luotettavampaa, ja sen avulla voidaan poistaa tarpeettomia if-lauseita. Jos mallista poistetaan mahdottomat tilat, niitä ei tarvitse edes testata \cite{impossiblebetter}. Seuraavaksi muutetaan virheellinen malli oikeaksi.

Kuvitellaan, että ollaan rakentamassa rajapintaa verkossa toimivalle tietovisa-alustalle, jossa tietovisa koostuu useasta monivalintakysymyksestä. Mallinnus voi alkaa näin.

\begin{code}
    \begin{minted}{typescript}
type Question = {
  question: string
  answers: Array<string>
  correctAnswerIndices: Array<number>
}

type Quiz = {
  description: string
  questions: Array<Question>
}
         
\end{minted}
    \caption{Mahdollinen lähestymistapa yksinkertaiselle tietovisan mallinnukselle. Malli koostuu visasta (Quiz) ja se kysymyksistä (Question)}
    \label{code:ts_first_quiz}
\end{code}

Yksinkertaisen tietovisan saisi toimimaan kyseisellä mallilla. Yksi kysymys tukee jopa montaa oikeaa vastausta.
Mallintamisessa on kuitenkin ongelma. Se valehtelee. Sillä voidaan kuvata tiloja, joiden pitäisi olla mahdottomia.

\textbf{Listat voivat olla tyhjiä.} Malli ei pidä huolta, että tietovisassa on oltava vähintään yksi kysymys. Malli ei pidä myöskään huolta siitä, että kysymyksillä tulisi olla vähintään yksi vastaus. Tämä johtuu siitä, että \mintinline{typescript}|Array|-tietorakenne voi olla tyhjä.

Jos mallissa käytetään listaa, mutta se ei kuitenkaan todellisuudessa saisi olla tyhjä, niin ohjelmoijan on tarkistettava koodissa manuaalisesti, että listan sisältö on läsnä.

Vastauksena ongelmaan on luoda tyyppi listalle, joka yksiselitteisesti ei voi olla tyhjä (Koodiesimerkki \ref{code:ts_non_empty}).

\begin{code}
    \begin{minted}{typescript}
type NonEmptyArray<T> = {
    first: T
    rest: Array<T>
}
    \end{minted}
    \caption{Mahdollinen lähestymistapa yksinkertaiselle tietovisan mallinnukselle. Tyypissä on käytetty tyyppimuuttujaa \texttt{T}. Tyyppimuuttuja tarkoittaa sitä, että sen sijalle voi laittaa minkä tahansa tyypin. Kyse on myös parametrisestä polymorfista (parametric polymorphism).}
    \label{code:ts_non_empty}
\end{code}

Tyypillä on kenttä \mintinline{javascript}|first|, jonka arvon on oltava olemassa. Ja kenttä \mintinline{javascript}|rest|, jossa säilytetään listan loput alkiot. Tuon kentän lista taas on tavallinen \mintinline{javascript}|Array|. Pelkällä tyypin muutoksella voidaan pitää huolta että listassa on 0..n sijasta 1..n alkiota.

\begin{code}
    \begin{minted}[highlightlines={3,4,9}]{typescript}
type Question = {
    question: string
    answers: NonEmptyArray<string> 
    correctAnswerIndices: Array<number>
}

type Quiz = {
    description: string
    questions: NonEmptyArray<Question> 
}       
    \end{minted}
    \caption{Vaihtoehtoinen lähestymistapa tietovisan mallintamiselle, jossa käytetään itsemääritettyä NonEmptyArray-tyyppiä}
    \label{code:ts_non_empty_quiz}
\end{code}

Nyt tietovisassa on oltava vähintään yksi kysymys (Koodiesimerkki \ref{code:ts_non_empty_quiz}), ja jokaisella kysymyksellä vähintään yksi vastaus. Kuitenkin oikeita vastauksia ei ole välttämättä ainuttakaan, sillä kenttä \mintinline{typescript}|correctAnswerIncides| on edelleen tyypiltään \mintinline{typescript}|Array<number>|.

Staattisella tyyppijärjestelmällä voi pitää huolen, että mallia seurataan. Tarkastuksia tyhjyydestä ei tarvitse tehdä if-lauseilla. Malli ei kuitenkaan ole vieläkään täydellinen.

\textbf{Kenttien arvot voivat olla mitä tahansa.} Mikään ei pidä huolta, että esimerkiksi kentän \mintinline{typescript}|correctAnswerIncides| indeksit ovat \glsdisp{correctness}{oikeellisia}. Myös kaikki, joissa tyyppinä on \mintinline{typescript}|string| voivat ottaa vastaan mitä tahansa merkkijonoja. Myös tyhjiä merkkijonoja.

Merkkijonojen validointi TypeScriptin tyyppijärjestelmässä on kuitenkin hieman vaivalloista, eikä välttämättä tarpeellista. Kenttä \mintinline{typescript}|correctAnswerIncides| on kuitenkin ajatuksena erittäin huono, sillä malli ei millään tavalla määrittele miten se liittyy kysymykseen tarkemmin. Mikään ei estä laittamasta täysin päättömiä numeroita \textquote{oikeiksi} indekseiksi. Korjauskeinona on paras keino: kentän poistaminen.

Kentän voi poistaa kokonaan, jos mallinnusta viedään pidemmälle. Voidaan sopia, että on olemassa kahdenlaisia vastauksia. Oikeita tai vääriä. Tämä on varsin helppo mallintaa (Koodiesimerkki \ref{code:ts_answer_types})

\begin{code}
    \begin{minted}{typescript}
type CorrectAnswer = {
    type: "correct"
    answer: string
}
type IncorrectAnswer = {
    type: "incorrect"
    answer: string
}
    \end{minted}
    \caption{Oikeille ja väärille vastauksille omat tyypit}
    \label{code:ts_answer_types}
\end{code}

Uusilla tyypeillä voidaan mallintaa kysymyksiä \glsdisp{correctness}{oikeellisemmin}. Voidaan myös määrittää, että kysymyksellä on oltava vähintään yksi oikea vastaus, ja nolla tai useampi väärä vastaus (Koodiesimerkki \ref{code:ts_question_with_explicit_answer_types}).

\begin{code}
    \begin{minted}[highlightlines={3,4}]{typescript}
type Question = {
    question: string
    correctAnswers: NonEmptyArray<CorrectAnswer>
    incorrectAnswers: Array<IncorrectAnswer>
}
    \end{minted}
    \caption{Kysymykseen voi tarkentaa millaisia vastauksia hyväksytään}
    \label{code:ts_question_with_explicit_answer_types}
\end{code}

Malli ei enää valehtele. Voi kuitenkin olla, että tulevaisuudessa on tarve luoda visalle enemmän ominaisuuksia. Voidaan haluta rajoittaa, että kysymyksessä voi valita vain yhden vaihtoehdon, tai usean. Voidaan esimerkiksi haluta antaa mahdollisuus antaa vastauksena vapaata tekstiä \cite{impossiblebetter}. Muita tarpeita voi olla loputtomasti.

\textbf{Voidaan ottaa käyttöön summatyyppejä.} Sen sijasta, että lisättäisiin tyyppeihin uusia kenttiä kuten \mintinline{javascript}|isMultipleAnswerQuestion| tai \mintinline{javascript}|isFreeForm|, voidaan puolestaan luoda jokaiselle fundamentaalisesti erilaiselle asialle kokonaan oma tyyppinsä, ja yhdistää ne yhdeksi summatyypiksi. Näin pidetään huolta, että jokaisella asialla voi olla vain kenttiä, jotka kuuluvat sille.

\begin{code}
    \begin{minted}[highlightlines={20}]{typescript}
type MultipleAnswerQuestion = {
    type: "multipleAnswerQuestion"
    question: string
    correctAnswers: NonEmptyArray<CorrectAnswer>
    incorrectAnswers: Array<IncorrectAnswer>
}

type SingleAnswerQuestion = {
    type: "singleAnswerQuestion"
    question: string
    correctAnswer: CorrectAnswer
    incorrectAnswers: Array<IncorrectAnswer>
}

type FreeFormQuestion = {
    type: "freeFormQuestion"
    question: string
}

type Question = MultipleAnswerQuestion | SingleAnswerQuestion | FreeFormQuestion
    \end{minted}
    \caption{Kysymysten mallintaminen summatyypillä. Monivalintakysymykselle, yksinkertaiselle kysymykselle ja vapaatekstikentälliselle kysymykselle on jokaiselle oma tyyppi. Tyypit on nostettu yhteen summatyyppiin (Question)}
    \label{code:ts_sum_type_nice}
\end{code}

Kysymysten \mintinline{javascript}|type|-kentän perusteella voidaan tarkistaa millaisesta kysymyksestä on todella kyse, jonka jälkeen voidaan suorittaa koodia, joka liittyy vain kyseisen kysymystyypin toimintaan.

Todellisuutta mallintavilla malleilla saadaan luotettavampaa koodia, ja poistetaan turhia if-lauseita. Myös jos mallista poistetaan mahdottomat tilat, niitä ei tarvitse edes testata.

\subsection{Virheet mukaan mallintamiseen}

Ohjelmissa tapahtuu virheitä. Vaikka pyrittäisiin malleihin, joista on poistettu mahdottomat tilat, niin ohjelmia ajaessa tullaan silti törmäämään virheisiin.

Miten on hyvä mallintaa tilanteita, kun virheitä tapahtuu?

Yleisesti imperatiivisissa paradigmoissa virhetilanteissa metodit ja funktiot palauttavat \texttt{null}-arvoja, virheitä tai poikkeuksia.
Funktionaalisessa ohjelmoinnissa voidaan käyttää samoja keinoja, mutta useimmin funktioiden suorituksen räjäyttämisen sijasta funktiot palauttavat summatyyppejä, tai monadeja, joissa voi olla sisällä joko oikea palautus- tai virhetyyppi.

TypeScriptissä yksinkertaisimmillaan tämän voi toteuttaa yhdellä omalla tyypillä (Koodiesimerkki \ref{code:ts_null_sum}).

\begin{code}
    \begin{minted}{typescript}
type Maybe<T> = T | null
const mightWorkMightNot = (): Maybe<string> => {
    if (Math.random() > 0.5) {
        return null
    }
    return "hello world"
}
    \end{minted}
    \caption{Mahdollisesti puuttuvan paluuarvon malli. Maybe voi sisältää jonkin tyypin tai arvon \texttt{null}}
    \label{code:ts_null_sum}
\end{code}

Deklaratiivisesti funktio kertoo, että se saattaa palauttaa merkkijonon, tai saattaa olla palauttamatta. Ehkä se on merkkijono, ehkä se ei ole.
Tyypin käyttäminen on hieman tönkköä, sillä jos funktion oikea palautusarvo on null, on mahdotonta sanoa oliko funktiokutsu onnistunut vai ei.

Tyypin voi viedä pidemmälle, jos arvot kääritään objekteihin. Vaihdetaan nimi Result-tyypiksi, että tyyppi voi kapseloida myös virheitä. (Koodiesimerkki \ref{code:ts_better_result_sum}.)

\begin{code}
    \begin{minted}{typescript}
type Success<T> = { value: T }
type Failure = { error: Error }
type Result<T> = Success<T> | Failure 

const mightWorkMightNot = (): Result<string> => {
    if (Math.random() > 0.5) {
        return { error: new Error("Sorry, bad failure.") }
    }
    return { value: "hello world" }
}
    \end{minted}
    \caption{Vaihtoehtoinen malli mahdollisesti epäonnistuvalle paluuarvolle. Result on joko Success, tai Failure. Molemmissa tapauksissa arvo on kuitenkin yhden kentän objekti, jonka ansiosta tyyppiin voi tallentaa \texttt{null}-arvon ilman että tietoa menetetään}
    \label{code:ts_better_result_sum}
\end{code}


Näin Result-tyyppi voi sisältää myös arvon \mintinline{typescript}|null|, sillä Result-objektin tarkan tyypin voi nyt päätellä ilman sitä. Nyt tarkan tyypin voi määrittää \mintinline{typescript}|value| tai \mintinline{typescript}|error| kentän läsnäolosta tai puutteesta. Näin rakennettuna tyypille on myös helppo rakentaa apufunktioita, joilla selvitetään, onko paluuarvo onnistunut, vai ei. Voi myös luoda apufunktiot, joilla luoda Result-objekteja. (Koodiesimerkki \ref{code:result_helpers}).

\begin{code}
    \begin{minted}{typescript}
const isFailure = <T>(x: Result<T>): x is Failure => 
    "error" in x
const isSuccess = <T>(x: Result<T>): x is Success<T> => 
    !isFailure(x)
const failureOf = (error: string): Failure => 
    ({ error: new Error(error) })
const successOf = <T>(value: T): Success<T> => 
    ({ value })
    \end{minted}
    \caption{Apufunktioita Result-tyypin tarkastamiseen ja luomiseen TypeScriptissä}
    \label{code:result_helpers}
\end{code}

Näillä voidaan rakentaa yhdistettyjä funktioita, joissa paluuarvon luonti ja palauttaminen on eksplisiittistä ja selkeää. Koodiesimerkissä \ref{code:result_pipe_example} käsitellään Result-objekteja yhdistetyssä funktiossa. Jokaisessa yhdisteen funktiossa tarkistetaan, onko aiempi tulos epäonnistunut. Epäonnistumisen sattuessa virhearvo päästetään suoraan putken läpi.

\begin{code}
    \begin{minted}{typescript}
const helloMaybe = (): Result<string> => {
    if (Math.random() > 0.5) {
        return failureOf("Sorry, bad failure.")
    }
    return successOf("hello world")
}

const addExclamation = (x: Result<string>): Result<string> => {
    if (isFailure(x)) return x
    return successOf(x.value + "!")
}

const toUpperCase = (x: Result<string>): Result<string> => {
    if (isFailure(x)) return x
    return successOf(x.value.toUpperCase())
}

pipe([
    helloMaybe,
    addExclamation,
    toUpperCase
])()
// Palauttaa { value : "HELLO WORLD!}
// tai { error: Error("Sorry, bad failure.") }
    \end{minted}
    \caption{Esimerkki yhdistetystä funktiosta Result-tyypin kanssa. Jokaisen välifunktion on tarkastettava onko edellinen kutsu onnistunut vai ei}
    \label{code:result_pipe_example}
\end{code}

Result-tyypin käyttö on selkeää, ja TypeScriptillä on pakko tarkistaa, onko arvo \mintinline{typescript}|Success| vai \mintinline{typescript}|Failure|, ennen kuin sen käärittyä arvoa voi tarkastella tai muokata.

Jos yksikin funktio palauttaa Result-objektin, jossa on virhe, objekti kulkee muuttumattomana loppujen funktioiden lävitse.

Kuitenkin ongelma on, että tämänkaltaisena Result-tyypin käyttäminen on loppujen lopuksi vain syntaksisokeroitu null-arvon tarkistus. Ne if-lauseet, joita aiemmin poistettiin todenmukaisella tyyppimallinnuksella, ovat tulleet takaisin aina pakollisen virhetarkistamisen muodossa.

Result-tyypillä on kuitenkin jotain, jolla sen voi viedä vielä askelta pidemmälle. Result-tyypin voi toteuttaa monadina.


\subsection{Result-monadi}

Monadi oli suuressa osassa teoriaosuudessa. Jatketaan mahdollisesti epäonnistuvan ohjelman toteuttamista siten, että siinä on mukana monadisia periaatteita. Monadia käyttämällä voidaan siirtää virheentarkastuksen if-lauseet yhteen paikkaan keskitetysti.

Luodaan \texttt{Result}-monadi (Koodiesimerkki \ref{code:result_monad_impl}).


\begin{code}
    \begin{minted}{typescript}
class Result {
    constructor(x) {
        this.value = x
    }
    static of(x) {
        return new Result(x)
    }
    bind(f) {
        if (this.value instanceof Error) {
            return this
        }

        const result = f(this.value)

        if (!(result instanceof Result)) {
            return Result.of(result)
        }

        return x
    }
}
\end{minted}
    \caption{Result implementoituna monadiksi. Result monadeja voi luoda staattisella \mintinline{js}|Result.of|-metodilla, ja monadiin voi ketjuttaa operaatioita \mintinline{js}|Result.prototype.bind|-metodilla. Jos ketjutuksessa monadin arvo on Error, niin ketjutus lopetetaan, ja saatu virhearvo kuljetetaan suoraan lävitse}
    \label{code:result_monad_impl}
\end{code}

Jos \texttt{Result}-monadi otetaan käyttöön aiemman osan koodiesimerkkiin, huomataan luettavan koodin määrän pienentyneen (Koodiesimerkki \ref{code:result_monad_example_2}).

\begin{code}
    \begin{minted}{typescript}
const helloMaybe = (): string | Error => 
    Math.random() > 0.5 
        ? "hello world" 
        : new Error("Sorry, bad failure.")

const addExclamation = (x: string) => x + "!"
const toUpperCase = (x: string) => x.toUpperCase()
const reverse = (x: string) => x.split("").reverse().join("")

Result.of(helloMaybe())
 .bind(addExclamation)
 .bind(toUpperCase)
 .bind(reverse)
\end{minted}
    \caption{Monadisia operaatioita ketjutettuna Result-monadiin. Annettavat funktiot voivat palauttaa käärimättömiä arvoja johtuen Result-monadin implementaatiosta}
    \label{code:result_monad_example_2}
\end{code}

Jos yksikin monadin funktioketjusta epäonnistuu, suoritus lakkaa ja virhe palautuu loppuun asti. Arvon voi vetää ulos monadista useilla eri tavoilla. Implementaation ansiosta tässä yhteydessä arvon voi ottaa ulos monadista yksinkertaisesti käyttäen \texttt{value}-kenttää (Koodiesimerkki \ref{code:result_monad_example_3})


\begin{code}
    \begin{minted}{typescript}
const monad = Result.of(helloMaybe())
    .bind(addExclamation)
    .bind(toUpperCase)
    .bind(reverse)
 const value = monad.value
 // value on joko "!DLROW OLLEH" tai Error("Sorry, bad failure")
\end{minted}
    \caption{Arvon voi poistaa monadin kontekstista yksinkertaisesti hakemalla \mintinline{js}|value|-kenttää }
    \label{code:result_monad_example_3}
\end{code}

Eri ohjelmointikielissä tai kirjastoissa monadeista arvon ulos saaminen voi poiketa. Matemaattisessa määritelmässä arvoa ei ole tarkoitus koskaan saada pois monadista. Esimerkiksi Haskellissa lähtökohtaisesti arvo kuuluu aina jättää monadinsa kontekstiin, eikä sitä tulisi käsitellä suoraan \cite{haskellallmonad}.

Result-monadin luominen ja käyttäminen vaatii ohjelmiston kehittäjiltä ylimääräistä kontekstia funktionaalisista käsitteistä. Kontekstin poistamiseksi voidaan käyttää räätälöidyn monadin sijasta JavaScriptin \texttt{Promise}-tietorakennetta (Koodiesimerkki \ref{code:result_monad_example_4}).

\begin{code}
    \begin{minted}{typescript}
const mightSayHello = async () => {
    if (Math.random() > 0.5) {
        throw new Error("Sorry, bad failure")
    }
    return 'hello world'
}
    
const promise = Promise.resolve(mightSayHello())
    .then(x => x+"!")
    .then(x => x.toUpperCase())
    .then(x => x.split("").reverse().join(""))
//  .catch(x => x)
\end{minted}
    \caption{Result-monadi-esimerkki käyttäen JavaScriptin sisäänrakennettua Promise-tietorakennetta. Toimintaperiaate vastaa aiempaa koodiesimerkkiä, eikä vaadi erillisiä kirjastoja}
    \label{code:result_monad_example_4}
\end{code}

\texttt{Promise}:n käyttäminen Result-monadin sijasta voi selkeyttää ohjelmakoodin luettavuutta JavaScript-kehittäjille. Toisaalta \texttt{Promise}:n käyttäminen esimerkin tavalla tuo turhaa monimutkaisuutta, sillä \texttt{Promise}:n luontainen asynkronisuus tuo turhia tehokkuushaittoja, kun operaatiot eivät muuten olisi asynkronisia \cite{mdn_promise}.

Ohjelmistokehittäjä Jake Archibald suosisi ohjelman kirjoittamista vielä enemmän JavaScriptin tyylillä (Koodiesimerkki \ref{code:archibald_bro}) \cite{pennane_fp_gist} .

\begin{code}
    \begin{minted}{typescript}
function someWellNamedFunction(): string | null {
    const result = mightSayHello();
    if (result === null) return;
    return reverse(result + '!').toUpperCase();
}
\end{minted}
    \caption{Ohjelmistokehittäjä Jake Archibaldin versio funktiosta \cite{pennane_fp_gist}. Versio suosii imperatiivisia tapoja toteuttaa funktio}
    \label{code:archibald_bro}
\end{code}

Saman asian voi siis toteuttaa monella eri tavalla. On suositeltavaa asettaa tiukan arvioinnin alle se, mikä versio sopii mihinkin tilanteeseen.



% Material and Methods
\clearpage%if the chapter heading starts close to bottom of the page, force a line break and remove the vertical vspace
\vspace{21.5pt}
\chapter{Kokemukset}

Tässä luvussa käsitellään ohjelmointikokemuksia erityisesti funktionaalisen ja oliopohjaisen ohjelmoinnin näkökulmista. Pohditaan, miten kokemukset, ohjelmointiparadigmat ja käytetyn kielen ominaisuudet vaikuttavat luettavuuteen, uudelleenkäytettävyyteen ja tehokkuuteen.

Kokemuksella on suuri rooli siinä, miten ohjelmoija lähestyy näitä teemoja. Vastaavasti eri ohjelmointiparadigmat — kuten funktionaalinen ja olio-ohjelmointi — tarjoavat erilaisia työkaluja ja periaatteita, jotka voivat joko helpottaa tai hankaloittaa näiden tavoitteiden saavuttamista. Ohjelmointikielen ominaisuudet, kuten Pythonin yksinkertaisuus tai JavaScriptin toiminnalliset työkalut, voivat myös muokata koodikäytäntöjä ja vaikuttaa siihen, mikä on optimaalista kussakin tilanteessa.

\section{Luettavuus}

Puidaan luettavuutta hieman eri näkökulmista. X-alustalla keskusteltiin entisen Google kehittäjän Jake Archibaldin kanssa, joka suostui vastaamaan funktionaaliseen ohjelmointiin liittyviin kysymyksiin \cite{pennane_x_convo}.

Yli 20 vuotta JavaScriptia ohjelmoinut Archibald pyrkii ohjelmoimaan koodia, joka on kuin suoraan ohjelmoinnin alkeita käsittelevästä oppikirjasta \cite{is_reduce_bad,pennane_fp_gist}. Hän pyrkii olemaan käyttämättä monimutkaisempia ohjelmoinnin työkaluja \cite{is_reduce_bad}, kuten esimerkiksi \texttt{Array.prototype.reduce}, osittaissoveltaminen (partial application) tai säännölliset lausekkeet (regular expressions) \cite{is_reduce_bad,pennane_x_convo}.

Hän suosii for-silmukoiden käyttämistä funktionaalisen \texttt{reduce}-funktion sijasta, ja perustelee valintaa luettavuudella \cite{is_reduce_bad,pennane_x_convo,pennane_fp_gist}. Yleisesti Archibald kuitenkin kokee funktionaalisen ohjelmoinnin olevan tuonut positiivisa vaikutteita JavaScriptiin \cite{pennane_x_convo}. Luettavuuden kannalta merkittävimpiä tekijöitä ovat keskimääräinen rivin pituus, tunnisteiden määrä per rivi ja keskimääräinen sulkeiden määrä \cite[8]{busereadability}. Edsger Dijkstran mukaan ohjelman luettavuus riippuu suurelta osin sen \glsdisp{ohjausrakenne}{ohjausrakenteiden} yksinkertaisuudesta \cite{dijkstra1976discipline}. Ohjelmoijat oppivat ensin for-silmukoita, ennen kuin ymmärtävät, mitä \texttt{reduce} (tai \texttt{fold}) tarkoittaa. Kuitenkin tunnisteet ja rakenteet, jotka ohjelmoija tuntee, ovat henkilökohtaisia ja perustuvat kokemuksiin sekä opiskelutaustaan.

Myös Python ohjelmointikielen luojana tunnettu Guido van Rossum halusi poistatti Python 3 versiosta \texttt{reduce}-funktion, ja perusteli sitä luettavuudella, ja kielen yksinkertaistamisella \cite{vanrossum2008reduce}.

Voi siis ajatella, että luettavuus on tarpeellista suhteuttaa työympäristöön. Jos koodikannassa suositaan funktioiden yhdistelemistä, tätä käytäntöä kannattaa noudattaa. Jos taas for-silmukoita käytetään mieluummin \texttt{reduce}-funktion sijasta, ne ovat parempi valinta.

Kokemuksien mukaan voi olla vaikeaa olla käyttämättä työkaluja, joiden näkee sopivan tilanteeseen. Kuitenkin kannattanee ottaa askel taaksepäin ja pitää huolta, että ohjelmakoodi on aina luettavaa tilanteeseen sopien.

\section{Uudelleenkäytettävyys}

Työpaikalla keskusteluissa on näkynyt käytänteitä, että ennen kuin jotakin tarvitaan kolmesti, ei sille ole kannattavaa kirjoittaa uudelleenkäytettävää abstraktiota.

Toisaalta kun kehitetään ulkoisia rajapintoja, on tilanne eri, sillä on pidettävä huolta taaksepäin suuntautuvasta yhteensopivuudesta.

Uudelleenkäytettävien palikoiden löytäminen ei myöskään ole itsestäänselvyys. On keskusteltu siitä, että kannattaisi välttää logiikan koostamista, ja toisinaan suosia yksinkertaisuuksia ja kovakoodattua tietoa, jotta sen etsiminen olisi koodikannasta helpompaa (esimerkiksi käyttämällä grep-komentorivikomentoa).

Uudelleenkäytettävyyttä voidaan harjoittaa erinäköisillä koodikäytänteillä, kuten erinäköisin funktionaalisen ohjelmoinnin periaattein, tai olio-ohjelmoinnin periaattein. Archibaldin mukaan koodikäytänteet usein epäonnistuvat tehtävässään, kun niissä yritetään liikaa pysytellä niiden puhtaissa akateemisissa raameissa \cite{pennane_fp_gist}.



\section{Tehokkuus}
Tehokkuutta ei ole suotavaa miettiä liikaa ennen kuin ohjelmoidaan systeemejä, jotka todella riippuvat siitä. Monet ovat tätä mieltä \cite{pennane_fp_gist,prematureoptimization}. Kun tehokkuusongelmat alkavat näkyä, on kokemusten mukaan todennäköisempää, että yhden ohjelman kriisipesäkkeen (hot spot) korjaamisella tilanne voidaan korjata, eikä tuhannen sivalluksen kuolema ole todennäköisin.

Ohjelman tehokkuuteen vaikuttavien asioiden tiedostaminen kuitenkin ennaltaehkäisee ongelmatilanteita. Esimerkiksi tiedon kopioinnin tai \texttt{reduce}-funktion välttämisellä voidaan ajaa proaktiivisesti ohjelmaa toimivampaan suuntaan.



% Material and Methods
%\clearpage%if the chapter heading starts close to bottom of the page, force a line break and remove the vertical vspace
\vspace{21.5pt}
\chapter{Johtopäätökset}

Funktionaalista ohjelmoinnin periaatteita voi ottaa mukaan perinteisempään ohjelmointiin. On kyseessä sitten funktionaalinen ohjelmointi tai olio-ohjelmointi, liiallinen pedanttisuus käsitteitä kohtaan johtaa ohjelman ylläpidettävyyden ongelmiin, kun tarkoituksena on tuottaa muuttuvaa koodia asiakaskäyttöön.

Asia jossa pedanttisuus on paikallaan on ongelmien mallintaminen. Kun ongelmat mallinnetaan oikein, säästytään turhalta työltä testauksen ja liiallisten if-lauseiden parissa. Oikein mallinnettuna ohjelmakoodin määrä pienenee, ja samalla ylläpidettävyys kasvaa.

Funktionaalisen ohjelmoinnin akateemiset käsitteet ovat tiukasti määriteltyjä. Mikään ei pakota seuraamaan määritelmiä pilkuntarkasti. Ei suosita sokeaa sääntöjen noudattamista ilman oikeaa perusteltua syytä. Yhdistetyt funktiot, funktioiden puhtaus ja tiedon muuttumattomuus, algebralliset tietotyypit ja rakenteet ovat toimivia työkaluja rakentaa ohjelmia, ja niitä voi hyödyntää mallin implemenointiin. Jos ohjelmointikieli ei tue niitä natiivisti, niiden mukaan tuominen kuitenkin sumuttaa koodin luettavuutta. Tärkeää on, että koodikannassa kunnioitetaan määritettyjä käytänteitä. Kokemusten mukaan keskinkertaisetkin käytänteet ovat parempia kuin se, ettei käytänteitä olisi lainkaan.

\section{Jatkotutkimus}

Funktionaalinen ohjelmointi on laaja ja monimuotoinen alue, joka kattaa teoreettisia käsitteitä ja käytännön sovelluksia. Jatkotutkimuksena voitaisiin tutkia useita eri näkökulmia, ohjelmontikieliä ja tunnettujen ohjelmoijien julkaisuja. Jatkotutkimus kannattaa pitää moniulotteisena, eikä pelkästään funktionaalisen ohjelmoinnin periaatteiden tutkiminen tuo välttämättä kykyä luoda pragmaattista ohjelmointitaitoa.


\subsection{Kategoriateoria}

Kategoriateorian jatkotutkiminen voi tarjota syvällisempää ymmärrystä funktionaalisen ohjelmoinnin periaatteista ja rakenteista. Kategoriateoria auttaa ymmärtämään funktioiden ja rakenteiden suhteita, mikä on keskeistä funktionaalisessa ohjelmoinnissa. Tämä voi johtaa tehokkaampiin ja elegantimpiin koodiratkaisuihin, joissa turha toisto poistuu.

Kun teoriaa ymmärtää, sitä ei ole silti pakko ottaa käytäntöön, kuitenkin kokemusten perusteella teorian läsnäolon tunnistaminen vahvistaa omia näkemyksiä siitä, miten toimia ja olla toimimatta.

\subsection{Fantasyland-spesifikaatio}

Fantasyland-spesifikaatio on määritelty joukko vaatimuksia ja sääntöjä, jotka koskevat funktioiden yhdistämistä ja algebrallisten rakenteiden käyttöä JavaScriptissä. Fantasyland tarjoaa mallin, joka perustuu kategoriateorian perusperiaatteisiin, erityisesti algebrallisiin rakenteisiin, kuten monadeihin ja funktioihin.

Fantasyland-spesifikaatio määrittelee, miten eri funktiot ja rakenteet voivat vuorovaikuttaa toistensa kanssa, tarjoten siten kehittäjille selkeät säännöt ja odotukset toiminnallisuuden implementoinnille. Tämä voi helpottaa erilaisten funktioiden ja datarakenteiden yhdistämistä siten, että niiden käytön johdonmukaisuus ja yhteensopivuus paranee.

Esimerkiksi, spesifikaatiossa määritellään monadeille ja niiden yhdistämiselle liittyvät säännöt, kuten \texttt{of} ja \texttt{bind} -operaatiot, jotka ovat keskeisiä funktioiden yhdistämisessä. Nämä säännöt perustuvat kategoriateorian käsitteisiin, mikä mahdollistaa funktioiden ja rakenteiden välisen yhteensopivuuden.

Fantasylandin periaatteiden ymmärtäminen voi auttaa kehittäjiä rakentamaan ohjelmia, joissa käytetään funktionaalisen ohjelmoinnin parhaita käytäntöjä. Tämä ei ainoastaan paranna koodin luettavuutta ja ylläpidettävyyttä, vaan myös mahdollistaa tehokkaamman koodin kehittämisen, joka hyödyntää algebrallisten rakenteiden tarjoamia etuja.

\subsection{Fp-ts-kirjasto}

Fp-ts-kirjaston tutkiminen auttaa puhtaan funktionaalisen ohjelmoinnin ymmärtämiseen TypeScript ympäristössä. Fp-ts-kirjasto on


\subsection{Eri ohjelmointikielet}

Eri ohjelmontikielien tutkiminen avaisi uusia näkökulmia funktionaalisen ohjelmoinnin käytäntöihin eri konteksteissa. Kielillä on omat erityispiirteensä, jotka rikastuttavat funktionaalisen ohjelmoinnin ymmärtämistä.

Esimerkiksi Haskell on yksi puhtaimmista funktionaalisista ohjelmointikielistä, ja sen tiukka tyyppijärjestelmä pakottaa kehittäjät ajattomaan ajatteluun ja puhtaiden funktioiden käyttöön \cite{haskellallmonad,haskellcomposition,haskellmonadlaws}. Haskellin kautta voisi vahvistaa ajatusta keskeisistä funktionaalisen ohjelmoinnin käsitteistä. Haskellin akateeminen luonne antaa syvällistä ymmärrystä funktionaalisen ohjelmoinnin teoreettisista perusteista.

Toisena Elixir on moderni kieli, joka yhdistää funktionaalisen ohjelmoinnin periaatteet ja ohjelmistojen rinnakkaisuuden \cite{elixir}. Se perustuu Erlangin ekosysteemiin, ja sen kyky käsitellä rinnakkaisia prosesseja auttaa opiskelijoita ymmärtämään, miten funktionaaliset ohjelmat voivat toimia tehokkaasti ja luotettavasti moniydinympäristöissä. Elixir korostaa myös ohjelmien luettavuutta ja ylläpidettävyyttä, mikä on tärkeä oppimistavoite funktionaalisessa ohjelmoinnissa.

Myös Go-ohjelmointikielen opiskelu voisi muuttaa merkittävästi ajatuksia miten ohjelmoida \cite{golang}. Go ei ole funktionaalinen ohjelmointikieli, vaikka se tukee funktionaalisia ohjelmointiperiaatteita, kuten korkeampia funktioita. Go:n käytännönläheinen lähestymistapa voi auttaa ymmärtämään, miten funktionaalisia periaatteita voidaan soveltaa (ja olla soveltamatta) ohjelmoinnissa.
\subsection{Tunnetut ohjelmistokehittäjät}

Funktionaalisen ohjelmoinnin ymmärtäminen vaatii monipuolista näkökulmaa eri lähteistä, sillä se on laaja ja monimuotoinen alue, joka ulottuu teoreettisista käsitteistä käytännön sovelluksiin. Eri ohjelmistokehittäjien työ tarjoaa arvokkaita esimerkkejä ja oivalluksia, jotka rikastuttavat ymmärrystä siitä, miten funktionaalisia periaatteita voidaan soveltaa ohjelmistokehityksessä.

Esimerkiksi Rich Hickeyn, Jake Archibaldin ja Richard Feldmanin töiden tarkastelu voisi avartaa näkemystä eri ohjelmointikielten erityispiirteitä ja niiden vaikutusta ohjelmistokehitykseen.

Rich Hickey on Clojure-ohjelmointikielen kehittäjä, ja hänen työnsä korostaa ohjelmoinnin yksinkertaisuutta ja tehokkuutta \cite{hickey_maybe_not,hickey_persistent_2009}. Hickeyn filosofia, joka perustuu datan muuttumattomuuteen ja puhtaisiin funktioihin, tarjoaa opiskelijoille syvällisen ymmärryksen siitä, miten ohjelmoinnin rakenteet voivat tukea luettavuutta ja ylläpidettävyyttä.

Jake Archibald on tunnettu erityisesti web-teknologioiden parissa työskentelevänä kehittäjänä \cite{against_self_closing_tags,is_reduce_bad}. Hänen näkemyksensä funktionaalisesta ohjelmoinnista, erityisesti JavaScriptin kontekstissa, auttaa opiskelijoita ymmärtämään, kuinka funktionaalisia periaatteita voidaan soveltaa käytännön projekteissa. Archibaldin työ tarjoaa käytännön esimerkkejä siitä, miten parantaa koodin laatua ja tehokkuutta web-kehityksessä.

Richard Feldman on tunnettu Elm-ohjelmointikielen puolestapuhuja, ja hänen työnsä keskittyy käyttöliittymäkehitykseen \cite{feldman_fp_pragmatists,impossiblebetter}. Feldmanin näkökulma korostaa funktionaalisen ohjelmoinnin etuja, kuten ennustettavuutta ja virheiden vähentämistä, mikä on erityisen hyödyllistä suurissa ohjelmistoprojekteissa. Hänen esimerkit auttavat opiskelijoita ymmärtämään, kuinka funktionaaliset periaatteet voivat parantaa sovellusten laatua ja käyttäjäkokemusta.

Yhteenvetona näiden kehittäjien työn tutkiminen tarjoaa käytännön esimerkkejä ja syvällisiä oivalluksia, jotka kasvattavat ymmärrystä  ohjelmoinnista.

\section{Loppusanat}

Työn alkuvaiheilla olin sitä mieltä että funktionaalinen ohjelmointi is epic big win. Nyt sitä mieltä, että funktionaalinen ohjelmointi on epäsuora voitto, eikä sitä kannata pakottaa jokaiseen tilanteeseen.

\clearpage%if the chapter heading starts close to bottom of the page, force a line break and remove the vertical vspace
\vspace{21.5pt}
\chapter{Yhteenveto}

Psykologi Kurt Lewin kiteytti aikoinaan, että mikään ei ole niin käytännöllistä kuin hyvä teoria. Tämä näkyy funktionaalisessa ohjelmoinnissa sekä hyvässä että pahassa: sen käsitteillä voidaan rakentaa teoreettisesti tehokkaita, luotettavia ja kompakteja ohjelmia, mutta sokea teorian seuraaminen voi johtaa tehokkuushaittoihin, luettavuusongelmiin ja liian nopeaan abstrahointiin.

Ohjelmien mallinnuksessa teoreettisuus ja funktionaalisen ohjelmoinnin deklaratiivisuus ovat hyödyksi. Mallit, jotka eivät valehtele vähentävät ylläpidettävän koodin määrää ja testaamisen tarvetta. Mallintamista voidaan viedä pidemmälle rakenteilla kuten monadeilla, jolloin mallin ja toteutuksen raja hämärtyy.

Työn alkuvaiheilla olin vakuuttunut funktionaalisen ohjelmoinnin merkittävästä
edusta. Nyt koen, että funktionaalinen ohjelmointi on enemmän kuin epäsuora voitto, eikä sitä tule pakottaa jokaiseen tilanteeseen. Kuvainnollisesti sanottuna kaikukammioon seinä on räjäytetty.

Ohjelmointia tehdään yhteistyössä, joten yhteisiä käytäntöjä on syytä noudattaa. Funktionaalisen ohjelmoinnin käsitteet ja käytänteet ovat ilmaisuvoimaisia ja ennustettavia, mutta myös erittäin teoreettisia ja tiukkoja, mikä vaatii opiskelua ja perustelua.

Löydettiin käytänteitä, joita on luonteva ottaa käyttöön tuleviin projekteihin, ja käytänteitä, joiden käyttöä kannattaa harkita tarkasti. Tutkimuksen lähteissä funktionaalisen ohjelmoinnin hyödyt ilmenivät parhaiten alunperin funktionaalisissa ohjelmointikielissä. Kunkin kielen ominaisuudet vaikuttavat merkittävästi siihen, kuinka tehokkaasti ja sujuvasti funktionaalisia periaatteita voidaan soveltaa.

%----------------------------------------------------------------------------------------
%	BIBLIOGRAPHY REFERENCES
%----------------------------------------------------------------------------------------

% Bibliography.
% Normally you do not edit this file.
% To add bibliography in your text, add them first in the biblio.bib file and
% reference them with the \cite{} command in your text. Check also \citeauthor
% and \citeyear (e.g. for copied figure caption), the \cites for multi-cites
% with page numbers,...

\clearpage
%line space
%\singlespacing %removed otherwise the appendix are also single space
\begin{singlespacing}
    \vspace*{-16pt}
    \printbibliography[title=\IfLanguageName{finnish}{Lähteet}{References}]
\end{singlespacing}

%for conting the pages
\label{LastPage}~


%----------------------------------------------------------------------------------------
%	APPENDICES
%----------------------------------------------------------------------------------------

% % Appendix
% Normally, you do not have to edit this file.

%start appendix
\appendix
%no page number for appendix in table of content
\addtocontents{toc}{\cftpagenumbersoff{chapter}}
%appendix sections and subsections not in table of content
\settocdepth{chapter}
%add "Appendices" in the table of content
\addappheadtotoc
%have Appendix 1 (instead of Appendix A)
\renewcommand{\thechapter}{\arabic{chapter}}

\newcommand\liite[1]{
    %each appendix start on a new page
    \clearpage
    %each appendix restart page num to one
    \setcounter{page}{1}
    %special counter for appendix TODO: this is a ugly quick hack :( Should find a better way to count the page per appendix.
    \newtotcounter{appx#1}
    %overwrite the header
    \makeevenhead{plain}{}{}{\appname \thechapter \\ \thepage\,(\stepcounter{appx#1}\total{appx#1})}
    \makeoddhead{plain}{}{}{\appname \thechapter \\ \thepage\,(\stepcounter{appx#1}\total{appx#1})}}



% %force smaller vertical spacing in table of content
% %!!! There can be some fun depending if the appendices have (sub)sections or not :D
% % You will have to play with these numbers and eventually add the \vspace line  before
% % some \chapter and force another number.
% % To add more fun, time to time the table of content get wrong after a build :(
% \addtocontents{toc}{\vspace{11pt}}
% \pretocmd{\chapter}{\addtocontents{toc}{\protect\vspace{-24pt}}}{}{}

% \liite{1}% This is a hack to have right page numbering for each appendix. Make sure to
% % use a unique number for each appendix.
% \input{sample/Xappendix1.tex}% Sample content to demonstrate appendix in LaTeX. You
% % are safe to delete this lines (and the next samples) once you refreshed your LaTeX
% % skills (and safe to delete the sample folder and all its file too).

% %\addtocontents{toc}{\vspace{11pt}}%fix vertical space for Table of Content
% \liite{2}
% \input{sample/Xappendix2.tex}

% \addtocontents{toc}{\vspace{11pt}}
% \liite{3}
% \input{sample/X_R_example.tex}


%----------------------------------------------------------------------------------------
%	THIS IS THE END
%----------------------------------------------------------------------------------------
\end{document}
